%% toc: P2 23 23.1,23.2,23.3,23.4
\section{Church's 1940 formulation of the simple theory of types}
\label{sec:p02-w02-p02-040-church}

The relational second-order language of
\cref{sec:second-order-logic-01} quantifies over relations of individuals
only, and the simple-type syntax of \cref{ch:simple-type-theory} fixes types
and terms without choosing axioms or rules.  This section gives one complete
typed higher-order calculus: the system of Church's paper of 1940
\cite{Church1940}.  It is a Hilbert-style calculus whose terms are typed
lambda-terms; comprehension is obtained from lambda-abstraction and
conversion, choice and descriptions from a selection constant, and
extensionality from an axiom schema.

\emph{Inputs.}  Finite strings, induction on formulas and on proofs, and, for the
semantic results, ZFC set theory as metatheory.
\emph{Conclusions.}  The exact primitive symbols, formation rules,
abbreviations, rules of inference and axioms of the system
(\cref{def:p02-w02-p02-040-church-types,def:p02-w02-p02-040-church-wff,def:p02-w02-p02-040-church-abbrev,def:p02-w02-p02-040-church-rules,def:p02-w02-p02-040-church-axioms});
the deduction theorem and the induction theorem; the quantifier, equality,
comprehension and arithmetic theorems selected below; soundness for the
standard interpretation, with consistency and the independence results
the source states; and a term-by-term bridge to the host simple-type syntax.
\emph{Standing.}  Notation in this section follows the source.  The
source's informal interpretation is made precise here as the standard model
of \cref{def:p02-w02-p02-040-church-standard}; that formalization is the
handbook's, and the source gives no model theory of its own.

\subsection{Types, symbols and well-formed formulas}

\begin{hbdefinition}{def:p02-w02-p02-040-church-types}{Church's type symbols}
The \emph{type symbols} form the least class of strings that contains
\(\iota\) and \(o\) and contains \((\alpha\beta)\) whenever it contains
\(\alpha\) and \(\beta\) \cite[\S1, p.~56]{Church1940}.  In the intended
reading \(o\) is the type of propositions, \(\iota\) the type of
individuals, and \((\alpha\beta)\) the type of one-argument functions whose
arguments are of type \(\beta\) and whose values are of type \(\alpha\)
\cite[\S1, pp.~56--57]{Church1940}.  Thus the \emph{value} type is written
first.  Functions of several arguments are curried: \(o\iota\iota\) is the
type of two-place propositional functions of individuals.

Parentheses in type symbols are omitted with association to the left, so
\(o\iota\) abbreviates \((o\iota)\), \(\iota\iota\iota\) abbreviates
\(((\iota\iota)\iota)\), and \(\iota(\iota\iota)\) abbreviates
\((\iota(\iota\iota))\).  Further,
\[
 \alpha'\ \text{abbreviates}\ ((\alpha\alpha)(\alpha\alpha)),\qquad
 \alpha''\ \text{abbreviates}\ ((\alpha'\alpha')(\alpha'\alpha')),
\]
and so on \cite[\S1, p.~56]{Church1940}.  Greek letters
\(\alpha,\beta,\gamma\) are metalinguistic variables over type symbols.
Type symbols occur in the formal language only as subscripts on variables
and constants.
\index{simple theory of types!Church's type symbols}
\end{hbdefinition}

\begin{hbdefinition}{def:p02-w02-p02-040-church-wff}{Primitive symbols and well-formed formulas}
The primitive symbols are
\[
 \lambda,\ (,\ ),\ N_{oo},\ A_{ooo},\ \Pi_{o(o\alpha)},\
 \iota_{\alpha(o\alpha)},\ a_\alpha,\ b_\alpha,\ \ldots,\ z_\alpha,\
 \bar a_\alpha,\ \bar b_\alpha,\ \ldots
\]
for every type symbol \(\alpha\) \cite[\S2, p.~57]{Church1940}.  The first
three are \emph{improper}; the rest are \emph{proper}.  The proper symbols
\(N_{oo}\), \(A_{ooo}\), \(\Pi_{o(o\alpha)}\) and \(\iota_{\alpha(o\alpha)}\)
are \emph{constants}; for each \(\alpha\) there is one \(\Pi_{o(o\alpha)}\)
and one \(\iota_{\alpha(o\alpha)}\).  All other proper symbols are
\emph{variables}, infinitely many of each type.  The letter \(\iota\)
therefore has two roles: the type symbol of individuals and the name of the
selection constant, distinguished by position (a type symbol is always a
subscript).

A \emph{formula} is any finite sequence of primitive symbols.  The
\emph{well-formed formulas} (wffs) and their types are given by three rules
\cite[\S2, p.~57]{Church1940}:
\begin{enumerate}
 \item a single proper symbol is a wff whose type is its subscript;
 \item if \(\mathbf x_\beta\) is a variable and \(\mathbf M_\alpha\) is a
       wff of type \(\alpha\), then \((\lambda\mathbf x_\beta\mathbf M_\alpha)\)
       is a wff of type \(\alpha\beta\);
 \item if \(\mathbf F_{\alpha\beta}\) and \(\mathbf A_\beta\) are wffs of
       types \(\alpha\beta\) and \(\beta\), then
       \((\mathbf F_{\alpha\beta}\mathbf A_\beta)\) is a wff of type \(\alpha\).
\end{enumerate}
The wffs are the least class closed under these rules, and each wff has a
unique type.  An occurrence of \(\mathbf x_\beta\) is \emph{bound} if it lies
in a well-formed part of the form \((\lambda\mathbf x_\beta\mathbf M)\), and
\emph{free} otherwise.  Bold capitals are syntactic variables for wffs and
bold lower-case letters syntactic variables for variables, the subscript
giving the type \cite[\S2, pp.~57--58]{Church1940}.

Parentheses are omitted where no ambiguity results, restoring them by
well-formedness and otherwise by association to the left; a type subscript
\(o\) may be omitted, so an italic letter without subscript has type \(o\)
\cite[\S2, p.~58]{Church1940}.
\index{simple theory of types!Church's well-formed formulas}
\end{hbdefinition}

\begin{hbdefinition}{def:p02-w02-p02-040-church-abbrev}{Church's abbreviations}
The following conventions of abbreviation are fixed, the arrow being read
``stands for'' \cite[\S2, p.~58]{Church1940}:
\[
\begin{array}{r@{\ \longrightarrow\ }l}
 [\sim\mathbf A_o] & N_{oo}\mathbf A_o,\\
 {}[\mathbf A_o\vee\mathbf B_o] & A_{ooo}\mathbf A_o\mathbf B_o,\\
 {}[\mathbf A_o\mathbf B_o] & [\sim[[\sim\mathbf A_o]\vee[\sim\mathbf B_o]]],\\
 {}[\mathbf A_o\supset\mathbf B_o] & [[\sim\mathbf A_o]\vee\mathbf B_o],\\
 {}[\mathbf A_o\equiv\mathbf B_o] & [[\mathbf A_o\supset\mathbf B_o][\mathbf B_o\supset\mathbf A_o]],\\
 {}[(\mathbf x_\alpha)\mathbf A_o] & \Pi_{o(o\alpha)}(\lambda\mathbf x_\alpha\mathbf A_o),\\
 {}[(\exists\mathbf x_\alpha)\mathbf A_o] & [\sim[(\mathbf x_\alpha)[\sim\mathbf A_o]]],\\
 {}[(\text{\rotatebox[origin=c]{180}{\(\iota\)}}\mathbf x_\alpha)\mathbf A_o] & \iota_{\alpha(o\alpha)}(\lambda\mathbf x_\alpha\mathbf A_o),\\
 Q_{o\alpha\alpha} & \lambda x_\alpha\lambda y_\alpha[(f_{o\alpha})[f_{o\alpha}x_\alpha\supset f_{o\alpha}y_\alpha]],\\
 {}[\mathbf A_\alpha=\mathbf B_\alpha] & Q_{o\alpha\alpha}\mathbf A_\alpha\mathbf B_\alpha,\\
 {}[\mathbf A_\alpha\neq\mathbf B_\alpha] & [\sim[\mathbf A_\alpha=\mathbf B_\alpha]],\\
 I_{\alpha\alpha} & \lambda x_\alpha x_\alpha,\\
 K_{\alpha\beta\alpha} & \lambda x_\alpha\lambda y_\beta x_\alpha,\\
 0_{\alpha'} & \lambda f_{\alpha\alpha}\lambda x_\alpha x_\alpha,\\
 1_{\alpha'} & \lambda f_{\alpha\alpha}\lambda x_\alpha(f_{\alpha\alpha}x_\alpha),\\
 2_{\alpha'} & \lambda f_{\alpha\alpha}\lambda x_\alpha(f_{\alpha\alpha}(f_{\alpha\alpha}x_\alpha)),\\
 3_{\alpha'} & \lambda f_{\alpha\alpha}\lambda x_\alpha(f_{\alpha\alpha}(f_{\alpha\alpha}(f_{\alpha\alpha}x_\alpha))),\ \text{etc.},\\
 S_{\alpha'\alpha'} & \lambda n_{\alpha'}\lambda f_{\alpha\alpha}\lambda x_\alpha(f_{\alpha\alpha}(n_{\alpha'}f_{\alpha\alpha}x_\alpha)),\\
 N_{o\alpha'} & \lambda n_{\alpha'}[(f_{o\alpha'})[f_{o\alpha'}0_{\alpha'}\supset[[(x_{\alpha'})[f_{o\alpha'}x_{\alpha'}\supset f_{o\alpha'}(S_{\alpha'\alpha'}x_{\alpha'})]]\supset f_{o\alpha'}n_{\alpha'}]]],\\
 \omega_{\alpha''\alpha'\alpha'} & \lambda y_{\alpha'}\lambda z_{\alpha'}\lambda f_{\alpha'\alpha'}\lambda g_{\alpha'}\lambda h_{\alpha\alpha}\lambda x_\alpha(y_{\alpha'}(f_{\alpha'\alpha'}g_{\alpha'}h_{\alpha\alpha})(z_{\alpha'}(g_{\alpha'}h_{\alpha\alpha})x_\alpha)),\\
 \langle\mathbf A_{\alpha'},\mathbf B_{\alpha'}\rangle & \omega_{\alpha''\alpha'\alpha'}\mathbf A_{\alpha'}\mathbf B_{\alpha'},\\
 P_{\alpha'\alpha'''} & \lambda n_{\alpha'''}(n_{\alpha'''}(\lambda p_{\alpha''}\langle S_{\alpha'\alpha'}(p_{\alpha''}(K_{\alpha'\alpha'\alpha'}I_{\alpha'})0_{\alpha'}),\\
 \multicolumn{2}{r}{p_{\alpha''}(K_{\alpha'\alpha'\alpha'}I_{\alpha'})0_{\alpha'}\rangle)\langle0_{\alpha'},0_{\alpha'}\rangle(K_{\alpha'\alpha'\alpha'}0_{\alpha'})I_{\alpha'}),}\\
 T_{\alpha''\alpha'} & \lambda x_{\alpha'}[(\text{\rotatebox[origin=c]{180}{\(\iota\)}}x_{\alpha''})[(N_{o\alpha''}x_{\alpha''})[x_{\alpha''}S_{\alpha'\alpha'}0_{\alpha'}=x_{\alpha'}]]],\\
 P_{\alpha'\alpha'} & \lambda x_{\alpha'}(P_{\alpha'\alpha'''}(T_{\alpha'''\alpha''}(T_{\alpha''\alpha'}x_{\alpha'}))).
\end{array}
\]
Square brackets introduced by these abbreviations may be omitted when no
ambiguity results.  A bold dot replacing an omitted opening bracket extends
its scope as far forward as well-formedness permits; otherwise brackets are
restored by association to the left, subject to a precedence among bracket
kinds: those of \([\sim\mathbf A_o]\) and \([\mathbf A_o\mathbf B_o]\) have the
lowest rank, then those of the quantifier, description, \(=\) and \(\neq\)
forms, then those of \(\vee\), and those of \(\supset\) and \(\equiv\) have
the highest rank; lower-ranked brackets receive the smaller scope
\cite[\S2, pp.~58--59]{Church1940}.  Every abbreviation is eliminable, so
statements about wffs apply to the formulas they abbreviate.  In
particular \(=\) is not primitive: \([\mathbf A_\alpha=\mathbf B_\alpha]\)
says that every propositional function true of \(\mathbf A_\alpha\) is true
of \(\mathbf B_\alpha\).
\index{simple theory of types!Church's abbreviations}
\end{hbdefinition}

The intended interpretation is given in words only
\cite[\S2, p.~59]{Church1940}: \(\lambda\) is abstraction, juxtaposition is
application, \(N_{oo}\) is negation, \(A_{ooo}\) disjunction,
\(\Pi_{o(o\alpha)}\) the universal quantifier taken as a propositional
function of propositional functions, and \(\iota_{\alpha(o\alpha)}\) a
selection operator.  In each type \(\alpha'\) the formulas
\(0_{\alpha'},1_{\alpha'},2_{\alpha'},\ldots\) are meant to denote the
natural numbers, \(S_{\alpha'\alpha'}\) a function on all of \(\alpha'\)
that acts as successor on natural numbers, and \(N_{o\alpha'}\) the property
of being a natural number of type \(\alpha'\).  The formula
\(T_{\alpha''\alpha'}\) carries a natural number of type \(\alpha'\) to the
corresponding one of type \(\alpha''\); \(P_{\alpha'\alpha'''}\), adapted from
Kleene, yields the predecessor of a number of type \(\alpha'''\) in type
\(\alpha'\); and \(P_{\alpha'\alpha'}\) is the predecessor within one type.
The source records that this same-type predecessor cannot be defined
without the selection operator, deriving this from the independence of the
axiom of infinity, which in its \S4 is Axioms~7 and~8 taken together
\cite[\S2, p.~59; \S4, p.~61]{Church1940}; the claim is reported without
proof (\cref{def:p02-w02-p02-041-church-predecessor}).  The source deliberately leaves the nature of
the types \(o\) and \(\iota\) open and admits other interpretations
\cite[\S1, p.~57]{Church1940}.

\subsection{Rules of inference and axioms}

\begin{hbdefinition}{def:p02-w02-p02-040-church-rules}{Church's rules of inference}
The rules of inference are the following six \cite[\S3, p.~60]{Church1940}.
\begin{enumerate}
 \item[I.] (Alphabetic change of a bound variable.)  If
   \(\mathbf x_\beta\) has no free occurrence in a part \(\mathbf M_\alpha\) and
   \(\mathbf y_\beta\) has no occurrence in it, that part may be replaced by
   the formula obtained from \(\mathbf M_\alpha\) by writing
   \(\mathbf y_\beta\) at every occurrence of \(\mathbf x_\beta\).
 \item[II.] (Reduction.)  A part
   \(((\lambda\mathbf x_\beta\mathbf M_\alpha)\mathbf N_\beta)\) may be replaced
   by the formula obtained from \(\mathbf M_\alpha\) by writing
   \(\mathbf N_\beta\) at every occurrence of \(\mathbf x_\beta\), on the
   condition that no bound variable of \(\mathbf M_\alpha\) is
   \(\mathbf x_\beta\) or a free variable of \(\mathbf N_\beta\).
 \item[III.] (Expansion.)  The converse of II: if \(\mathbf A_\alpha\) is
   obtained from \(\mathbf M_\alpha\) by writing \(\mathbf N_\beta\) at every
   occurrence of \(\mathbf x_\beta\), a part \(\mathbf A_\alpha\) may be
   replaced by \(((\lambda\mathbf x_\beta\mathbf M_\alpha)\mathbf N_\beta)\),
   on the same condition.
 \item[IV.] From \(\mathbf F_{o\alpha}\mathbf x_\alpha\) infer
   \(\mathbf F_{o\alpha}\mathbf A_\alpha\), provided \(\mathbf x_\alpha\) is not
   a free variable of \(\mathbf F_{o\alpha}\).
 \item[V.] From \(\mathbf A_o\supset\mathbf B_o\) and \(\mathbf A_o\) infer
   \(\mathbf B_o\).
 \item[VI.] From \(\mathbf F_{o\alpha}\mathbf x_\alpha\) infer
   \(\Pi_{o(o\alpha)}\mathbf F_{o\alpha}\), provided \(\mathbf x_\alpha\) is
   not a free variable of \(\mathbf F_{o\alpha}\).
\end{enumerate}
A \emph{part} of a formula is a consecutive well-formed part other than a
variable immediately following an occurrence of \(\lambda\); when a part
occurs several times, any one occurrence may be replaced; and substituting
for a variable that does not occur leaves the formula unchanged
\cite[\S3, p.~60]{Church1940}.  Rules I--III are the rules of
\emph{\(\lambda\)-conversion}; a chain of their applications is a
\emph{conversion}.  Rule IV is \emph{substitution}, V \emph{modus ponens},
and VI \emph{generalization}.  The source states two derived rules
\cite[\S3, p.~60]{Church1940}:
\begin{enumerate}
 \item[IV\('\).] From \(\mathbf M_o\), infer the formula obtained by putting
   \(\mathbf A_\alpha\) in place of each free occurrence of
   \(\mathbf x_\alpha\), on the condition that no bound variable of
   \(\mathbf M_o\) other than \(\mathbf x_\alpha\) is free in
   \(\mathbf A_\alpha\).
 \item[VI\('\).] From \(\mathbf M_o\) infer \((\mathbf x_\alpha)\mathbf M_o\).
\end{enumerate}
The source obtains IV\('\) from conversion and IV, and VI\('\) from
conversion and VI.  Here is that reduction.  For IV\('\), choose by Rule~I a
variant of \(\mathbf M_o\) whose bound variables avoid \(\mathbf x_\alpha\);
Rule~III turns \(\mathbf M_o\) into
\((\lambda\mathbf x_\alpha\mathbf M_o)\mathbf x_\alpha\), in which
\(\mathbf x_\alpha\) is not free in the functor; IV gives
\((\lambda\mathbf x_\alpha\mathbf M_o)\mathbf A_\alpha\); and II, whose
proviso is the stated one, gives the substitution instance.  For VI\('\),
Rule~III gives \((\lambda\mathbf x_\alpha\mathbf M_o)\mathbf x_\alpha\), VI
gives \(\Pi_{o(o\alpha)}(\lambda\mathbf x_\alpha\mathbf M_o)\), which is
\((\mathbf x_\alpha)\mathbf M_o\).
\index{simple theory of types!Church's rules of inference}
\index{lambda-conversion!in Church's type theory}
\end{hbdefinition}

\begin{hbdefinition}{def:p02-w02-p02-040-church-axioms}{Church's formal axioms}
The formal axioms are all formulas of the following infinite list, one
instance of each superscripted schema for every choice of the type symbols
\(\alpha,\beta\) \cite[\S4, pp.~60--61]{Church1940}:
\[
\begin{array}{rl}
 1. & p\vee p\supset p.\\
 2. & p\supset p\vee q.\\
 3. & p\vee q\supset q\vee p.\\
 4. & p\supset q\supset.\,r\vee p\supset r\vee q.\\
 5^\alpha. & \Pi_{o(o\alpha)}f_{o\alpha}\supset f_{o\alpha}x_\alpha.\\
 6^\alpha. & (x_\alpha)[p\vee f_{o\alpha}x_\alpha]\supset p\vee\Pi_{o(o\alpha)}f_{o\alpha}.\\
 7. & (\exists x_\iota)(\exists y_\iota)\,.\,x_\iota\neq y_\iota.\\
 8. & N_{o\iota'}x_{\iota'}\supset.\,N_{o\iota'}y_{\iota'}\supset.\,S_{\iota'\iota'}x_{\iota'}=S_{\iota'\iota'}y_{\iota'}\supset x_{\iota'}=y_{\iota'}.\\
 9^\alpha. & f_{o\alpha}x_\alpha\supset.\,(y_\alpha)[f_{o\alpha}y_\alpha\supset x_\alpha=y_\alpha]\supset f_{o\alpha}(\iota_{\alpha(o\alpha)}f_{o\alpha}).\\
 10^{\alpha\beta}. & (x_\beta)[f_{\alpha\beta}x_\beta=g_{\alpha\beta}x_\beta]\supset f_{\alpha\beta}=g_{\alpha\beta}.\\
 11^\alpha. & f_{o\alpha}x_\alpha\supset f_{o\alpha}(\iota_{\alpha(o\alpha)}f_{o\alpha}).
\end{array}
\]
The \emph{theorems} are the formulas obtainable from the axioms by finitely
many applications of the rules; a \emph{proof} is a finite sequence each of
whose members is an axiom or follows from earlier members by a rule
\cite[\S4, p.~61]{Church1940}.  A \emph{variant} of an axiom is a formula
obtained from it by alphabetic changes of free and bound variables that
keep types, keep equal variables equal and different variables different;
variants are theorems by Rules I and IV\('\) \cite[\S5, p.~62]{Church1940}.

A \emph{proof of \(\mathbf B_o\) on the assumptions
\(\mathbf A^1_o,\ldots,\mathbf A^n_o\)} is a finite sequence ending in
\(\mathbf B_o\) each member of which is an assumption, a variant of an axiom,
or obtained from earlier members by a rule, subject to the condition that
no variable free in any assumption is substituted for (Rule IV) or
generalized upon (Rule VI).  Its existence is written
\(\mathbf A^1_o,\ldots,\mathbf A^n_o\derives\mathbf B_o\); for \(n=0\),
\(\derives\mathbf B_o\) says that \(\mathbf B_o\) is a theorem
\cite[\S5, p.~62]{Church1940}.
\index{simple theory of types!Church's axioms}
\index{axiom of choice!Church's type theory}
\index{extensionality!Church's type theory}
\index{descriptions!Church's type theory}
\end{hbdefinition}

The source assigns the axioms their roles \cite[\S4, pp.~61--62]{Church1940}:
Axioms~1--4 suffice for the propositional calculus, and Axioms~1--6\(^\alpha\)
for the logical functional calculus; Axioms~7 and~8 together act as an axiom
of infinity, and with 9\(^\alpha\) (axioms of \emph{descriptions}) they yield
elementary number theory; 10\(^{\alpha\beta}\) (\emph{extensionality} for
functions) and 11\(^\alpha\) (\emph{choice}) are needed for classical
analysis.  The source qualifies every such necessity claim in its n.~8:
contextual definitions, such as Russell's treatment of classes and
descriptions, are deliberately not used, and ``necessary'' is to be read
relative to that exclusion \cite[\S4, p.~61, n.~8]{Church1940}.

Descriptions are kept although 9\(^\alpha\) follows from Axioms 1--4 and
11\(^\alpha\) (\cref{thm:p02-w02-p02-040-church-9-from-11}), so that the
consequences of 1--9\(^\alpha\) without 10 and 11 can be studied
\cite[\S4, pp.~61--62]{Church1940}.  The source states that when 1--9\(^\alpha\)
are the only axioms, each 9\(^\alpha\) is independent
\cite[\S4, p.~62]{Church1940}; \cref{thm:p02-w02-p02-040-church-selection}
proves this.  It also observes that once 10\(^{\alpha\beta}\) is present,
only \(\iota_{o(oo)}\) and \(\iota_{\iota(o\iota)}\) need be primitive, the
higher selection constants being introduced as abbreviations by the schema
\[
 \iota_{\alpha\beta(o(\alpha\beta))}\longrightarrow
 \lambda h_{o(\alpha\beta)}\lambda x_\beta
 (\text{\rotatebox[origin=c]{180}{\(\iota\)}}y_\alpha)(\exists f_{\alpha\beta})\,.\,
 h_{o(\alpha\beta)}f_{\alpha\beta}\,.\,y_\alpha=f_{\alpha\beta}x_\beta,
\]
in such a way that the formulas 9\(^\alpha\), read through these
abbreviations, become theorems of Axioms 1--8, 9\(^o\), 9\(^\iota\) and
10\(^{\alpha\beta}\) \cite[\S4, p.~62]{Church1940}.  The source gives no proof;
\cref{thm:p02-w02-p02-040-church-selection} supplies one.

Axiom 10\(^{\alpha\beta}\) is weaker in some directions than other
extensionality axioms, and the source argues that it is adequate because a
class can be represented by a function into the numerals of type
\(\iota'\): the class associated with the propositional function denoted by
\(\mathbf F_{o\alpha}\) is denoted by
\[
 \lambda x_\alpha(\text{\rotatebox[origin=c]{180}{\(\iota\)}}y_{\iota'})\,.\,
 (\mathbf F_{o\alpha}x_\alpha)[y_{\iota'}=0_{\iota'}]\vee
 (\sim\mathbf F_{o\alpha}x_\alpha)[y_{\iota'}=1_{\iota'}]
\]
\cite[\S4, p.~61]{Church1940}, the juxtapositions inside the scope of the
description being conjunctions.  The source also considers the further
extensionality axiom \(p\equiv q\supset p=q\): it would leave only two
propositions and, together with 10\(^{\alpha\beta}\), would allow classes to
be identified with propositional functions \cite[\S4, p.~61]{Church1940}.

\begin{hbconvention}{def:p02-w02-p02-040-church-prop}{The propositional input (P)}
The source does not derive the theorems of the propositional calculus from
Axioms~1--4: it states that their derivation by Rules IV\('\) and V is
known and refers to Hilbert and Ackermann and to Bernays
\cite[\S5, p.~62 and n.~9]{Church1940}.  The derivations below use the
following fact, called \emph{(P)}: every substitution instance, by
wffs of type \(o\), of a classical propositional tautology in
\(\sim,\vee\) is a theorem.  The source's reference is recorded here; the
handbook does not rely on it, because
\cref{thm:p02-w02-p02-040-church-prop-complete} proves (P) from Axioms~1--4
by Rules I, IV\('\) and V.
\end{hbconvention}

\begin{hbtheorem}{thm:p02-w02-p02-040-church-prop-complete}{proved=annex:p02-w02-p02-040-church-prop-complete}{Axioms 1--4 yield every tautology instance}
Let \(\tau(p_1,\ldots,p_k)\) be a formula built from distinct variables
\(p_1,\ldots,p_k\) of type \(o\) by \(\sim\) and \(\vee\) alone (with
\(\supset\) read as its abbreviation), and suppose \(\tau\) takes the value
true under every assignment of truth values to \(p_1,\ldots,p_k\).  Then for
all wffs \(\mathbf B^1_o,\ldots,\mathbf B^k_o\),
\(\derives\tau(\mathbf B^1_o,\ldots,\mathbf B^k_o)\), by a proof whose axioms
are variants of Axioms~1--4 and whose rules are I, IV\('\) and V.
\end{hbtheorem}

\subsection{Selected derivations}

The source states theorems 12\(^\alpha\)--20\(^\alpha\) and
22\(^\alpha\)--24\(^\alpha\) as consequences of Axioms 1--6\(^\alpha\) and
leaves their proofs to the reader \cite[\S5, p.~63; \S6, p.~64]{Church1940};
the annex supplies them.

\begin{hbtheorem}{thm:p02-w02-p02-040-church-deduction}{proved=annex:p02-w02-p02-040-church-deduction}{Deduction theorem (Church's VII)}
If \(\mathbf A^1_o,\ldots,\mathbf A^n_o\derives\mathbf B_o\), then
\(\mathbf A^1_o,\ldots,\mathbf A^{n-1}_o\derives\mathbf A^n_o\supset\mathbf B_o\),
for \(n=1,2,3,\ldots\) \cite[\S5, pp.~62--63]{Church1940}.
\end{hbtheorem}

\begin{hblemma}{thm:p02-w02-p02-040-church-12-21}{proved=annex:p02-w02-p02-040-church-12-21}{Quantifier and equality theorems 12--21}
For every type \(\alpha\) (and \(\beta\) where it occurs) the following are
theorems; 12--20 use only Axioms 1--6\(^\alpha\), and 21 uses only
Axioms 1--4 and 10\(^{\alpha\beta}\) \cite[\S5, p.~63]{Church1940}:
\[
\begin{array}{rl}
 12^\alpha. & (x_\alpha)f_{o\alpha}x_\alpha\supset f_{o\alpha}y_\alpha.\\
 13^\alpha. & f_{o\alpha}y_\alpha\supset(\exists x_\alpha)f_{o\alpha}x_\alpha.\\
 14^\alpha. & (x_\alpha)[p\supset f_{o\alpha}x_\alpha]\supset.\,p\supset(x_\alpha)f_{o\alpha}x_\alpha.\\
 15^\alpha. & (x_\alpha)[f_{o\alpha}x_\alpha\supset p]\supset.\,(\exists x_\alpha)f_{o\alpha}x_\alpha\supset p.\\
 16^\alpha. & x_\alpha=x_\alpha.\\
 17^\alpha. & x_\alpha=y_\alpha\supset.\,f_{o\alpha}x_\alpha\supset f_{o\alpha}y_\alpha.\\
 18^{\beta\alpha}. & x_\alpha=y_\alpha\supset f_{\beta\alpha}x_\alpha=f_{\beta\alpha}y_\alpha.\\
 19^\alpha. & x_\alpha=y_\alpha\supset y_\alpha=x_\alpha.\\
 20^\alpha. & x_\alpha=y_\alpha\supset.\,y_\alpha=z_\alpha\supset x_\alpha=z_\alpha.\\
 21^{\alpha\beta}. & f_{\alpha\beta}=\lambda x_\beta(f_{\alpha\beta}x_\beta).
\end{array}
\]
\end{hblemma}

Theorem 21\(^{\alpha\beta}\) is the \(\eta\)-equation.  It is not a rule of
conversion in this system: Rules I--III are \(\alpha\)- and
\(\beta\)-conversion only, and the \(\eta\)-equation is obtained as a
theorem from extensionality.

\begin{hbtheorem}{thm:p02-w02-p02-040-church-comprehension}{proved=annex:p02-w02-p02-040-church-comprehension}{Comprehension by abstraction}
Let \(\mathbf A_o\) be any wff and \(\mathbf f_{o\alpha}\) a variable not
free in \(\mathbf A_o\).  Then, using only Axioms 1--6\(^\alpha\) and the
rules,
\[
 \derives(\exists\mathbf f_{o\alpha})(\mathbf x_\alpha)
 [\mathbf f_{o\alpha}\mathbf x_\alpha\equiv\mathbf A_o].
\]
The witness is the abstract \(\lambda\mathbf x_\alpha\mathbf A_o\), and
\(\mathbf A_o\) may contain quantifiers of any type, so the comprehension
obtained is impredicative.
\end{hbtheorem}

This theorem is not numbered in the source; it records how the source's
system supplies the comprehension that the relational calculus of
\cref{def:p02-sol-derived-calculus} postulates as the schema family \(\mathsf{Comp}^n_\chi\) of \cref{def:p02-sol-matched-henkin}.  The source likewise introduces
classes through abstracts, representing the class of \(F_{o\alpha}\) by an
abstract with values among the numerals \(0_{\iota'}\) and \(1_{\iota'}\)
\cite[\S4, p.~61]{Church1940}.

\begin{hbtheorem}{thm:p02-w02-p02-040-church-induction}{proved=annex:p02-w02-p02-040-church-induction}{Peano theorems 22--24 and the induction theorem (Church's VIII)}
For every type \(\alpha\), using only Axioms 1--6\(^\alpha\),
\[
\begin{array}{rl}
 22^\alpha. & N_{o\alpha'}0_{\alpha'}.\\
 23^\alpha. & N_{o\alpha'}x_{\alpha'}\supset N_{o\alpha'}(S_{\alpha'\alpha'}x_{\alpha'}).\\
 24^\alpha. & f_{o\alpha'}0_{\alpha'}\supset.\,(x_{\alpha'})[N_{o\alpha'}x_{\alpha'}\supset.\,f_{o\alpha'}x_{\alpha'}\supset f_{o\alpha'}(S_{\alpha'\alpha'}x_{\alpha'})]\supset.\,N_{o\alpha'}x_{\alpha'}\supset f_{o\alpha'}x_{\alpha'}
\end{array}
\]
are theorems \cite[\S6, p.~64]{Church1940}.  Moreover (VIII): if
\(\mathbf x_{\alpha'}\) is not free in
\(\mathbf A^1_o,\ldots,\mathbf A^n_o,\mathbf F_{o\alpha'}\), if
\(\mathbf A^1_o,\ldots,\mathbf A^n_o\derives\mathbf F_{o\alpha'}0_{\alpha'}\), and if
\(\mathbf A^1_o,\ldots,\mathbf A^n_o,N_{o\alpha'}\mathbf x_{\alpha'},
\mathbf F_{o\alpha'}\mathbf x_{\alpha'}\derives
\mathbf F_{o\alpha'}(S_{\alpha'\alpha'}\mathbf x_{\alpha'})\), then
\(\mathbf A^1_o,\ldots,\mathbf A^n_o\derives
N_{o\alpha'}\mathbf x_{\alpha'}\supset\mathbf F_{o\alpha'}\mathbf x_{\alpha'}\)
\cite[\S6, p.~64]{Church1940}.
\end{hbtheorem}

\begin{hbtheorem}{thm:p02-w02-p02-040-church-27-29}{proved=annex:p02-w02-p02-040-church-27-29}{Theorems 27, 28, 25 and 29}
For every type \(\alpha\) the following are theorems
\cite[\S6, pp.~64--65]{Church1940}:
\[
\begin{array}{rl}
 27^\alpha. & (\exists x_\alpha)(\exists y_\alpha)\,.\,x_\alpha\neq y_\alpha.\\
 28^\alpha. & S_{\alpha'\alpha'}x_{\alpha'}\neq0_{\alpha'}.\\
 25^\alpha. & N_{o\alpha'}x_{\alpha'}\supset S_{\alpha'\alpha'}x_{\alpha'}\neq0_{\alpha'}.\\
 29^\alpha. & N_{o\alpha''}n_{\alpha''}\supset N_{o\alpha'}(n_{\alpha''}S_{\alpha'\alpha'}0_{\alpha'}).
\end{array}
\]
Axiom~7 enters only as the case 27\(^\iota\); every other case uses only
Axioms 1--6\(^\alpha\) together with the earlier cases.
\end{hbtheorem}

Theorem 27\(^\alpha\) is proved by induction on types: 27\(^o\) is derived
from the propositional calculus with 17\(^o\) and 13\(^o\), 27\(^\iota\)
is Axiom~7, and 27\(^{\alpha\beta}\) is derived from 27\(^\alpha\).
The source's own route, followed in the annex, derives 28\(^\alpha\) from
27\(^\alpha\) through the constant functions \(K_{\alpha\alpha\alpha}z_\alpha\)
\cite[\S6, p.~65]{Church1940}.

\subsection{The successor axiom at higher types: theorems 26 and 30}

Theorem 26\(^\alpha\) (displayed in \cref{thm:p02-w02-p02-040-church-soundness})
is the Peano postulate that successor is one-to-one on the natural numbers
of type \(\alpha'\).  The source states that 26\(^\iota\) is Axiom~8 and that
26\(^{\iota'}\), 26\(^{\iota''}\), \ldots{} are theorems, ``consequences of
1--6\(^\alpha\) and 8'' \cite[\S6, p.~64]{Church1940}, and it proves this
through a companion formula \cite[\S6, p.~66]{Church1940}:
\[
 30^\alpha.\quad N_{o\alpha''}m_{\alpha''}\supset.\,N_{o\alpha''}n_{\alpha''}\supset.\,
 m_{\alpha''}S_{\alpha'\alpha'}0_{\alpha'}=n_{\alpha''}S_{\alpha'\alpha'}0_{\alpha'}
 \supset m_{\alpha''}=n_{\alpha''}.
\]
Read through \(T_{\alpha''\alpha'}\), the formula
\(m_{\alpha''}S_{\alpha'\alpha'}0_{\alpha'}\) is the number of type \(\alpha'\)
corresponding to the number \(m_{\alpha''}\), so 30\(^\alpha\) says that this
correspondence is one-to-one.  As with 26\(^\alpha\), not every 30\(^\alpha\)
is a theorem \cite[\S6, p.~66]{Church1940};
\cref{thm:p02-w02-p02-041-church-30} proves that 30\(^o\) is not a theorem.

\begin{hbtheorem}{thm:p02-w02-p02-040-church-26-30}{proved=annex:p02-w02-p02-040-church-26-30}{Theorems 26 and 30 for the types \(\iota,\iota',\iota'',\ldots\)}
For every type symbol \(\alpha\) \cite[\S6, pp.~66--67]{Church1940}:
\begin{enumerate}
 \item if \(\derives26^\alpha\) then \(\derives30^\alpha\);
 \item if \(\derives26^\alpha\) and \(\derives30^\alpha\) then
   \(\derives26^{\alpha'}\).
\end{enumerate}
Consequently 26\(^\alpha\) and 30\(^\alpha\) are theorems for
\(\alpha=\iota,\iota',\iota'',\ldots\).  The derivations written in the annex
use Axioms~1--8 only.
\end{hbtheorem}

The source's footnote 12 attaches three qualifications to these results
\cite[\S6, p.~66, n.~12]{Church1940}.  First, it raises the question whether
30\(^\iota\) could replace Axiom~8 as the second half of the axiom of infinity.
Second, it reports two partial answers, neither proved in the paper: the
author has a proof, using properties of \(P_{\iota'\iota'''}\), that 30\(^\iota\)
and 30\(^{\iota'}\) together can replace Axiom~8 in the presence of
1--6\(^\alpha\); and A.~M.~Turing has a proof that, in the presence of 1--7 and
9\(^\alpha\), 30\(^\iota\) alone can replace it.  Third, whether Axiom~8 is
independent of 1--7 together with 30\(^\iota\) is recorded as open, the usual
methods of eliminating descriptions not applying.  The handbook proves none
of these three statements.

\begin{hbwarning}{def:p02-w02-p02-040-church-26-axiom-7}{Which axioms the derivation of 26 uses}
The source's proof of clause~1 uses 28\(^\alpha\)
\cite[\S6, pp.~66--67]{Church1940}, and for \(\alpha=\iota,\iota',\ldots\)
the derivation of 28\(^\alpha\) in \cref{thm:p02-w02-p02-040-church-27-29}
passes through 27\(^\iota\), which is Axiom~7.  So the derivation of
26\(^{\iota'}\), 26\(^{\iota''}\), \ldots{} written here, which follows the
source's, uses Axioms~1--8, not only 1--6\(^\alpha\) and 8 as the source's
announcement on p.~64 says.  Whether Axiom~7 can be removed from it is not
settled here.  Axiom~7 is not removable by an interpretation with one
individual: there every type has one member, so 26\(^\alpha\) holds at every
type.
\end{hbwarning}

\subsection{Properties of \(T_{\alpha''\alpha'}\): theorems 31--35}

The formula \(T_{\alpha''\alpha'}\) of \cref{def:p02-w02-p02-040-church-abbrev}
is defined by a description: \(T_{\alpha''\alpha'}x_{\alpha'}\) is the number
\(x_{\alpha''}\) of type \(\alpha''\) with
\(x_{\alpha''}S_{\alpha'\alpha'}0_{\alpha'}=x_{\alpha'}\).  Its expected
properties need Axiom 9\(^{\alpha''}\) and, for uniqueness, 30\(^\alpha\);
the existence half needs neither \cite[\S7, p.~67]{Church1940}.

\begin{hbtheorem}{thm:p02-w02-p02-040-church-31-35}{proved=annex:p02-w02-p02-040-church-31-35}{Theorems 31--35 on \(T_{\alpha''\alpha'}\)}
For every type \(\alpha\), using Axioms 1--6\(^\alpha\) only,
\[
 33^\alpha.\quad N_{o\alpha'}x_{\alpha'}\supset(\exists x_{\alpha''})\,.\,
 N_{o\alpha''}x_{\alpha''}\,.\,x_{\alpha''}S_{\alpha'\alpha'}0_{\alpha'}=x_{\alpha'}
\]
is a theorem.  If \(\alpha\) is a type for which \(\derives30^\alpha\), then,
using in addition Axiom 9\(^{\alpha''}\), the following are theorems
\cite[\S7, pp.~67--68]{Church1940}:
\[
\begin{array}{rl}
 31^\alpha. & N_{o\alpha'}x_{\alpha'}\supset N_{o\alpha''}(T_{\alpha''\alpha'}x_{\alpha'}).\\
 32^\alpha. & N_{o\alpha'}x_{\alpha'}\supset T_{\alpha''\alpha'}x_{\alpha'}S_{\alpha'\alpha'}0_{\alpha'}=x_{\alpha'}.\\
 34^\alpha. & N_{o\alpha'}x_{\alpha'}\supset T_{\alpha''\alpha'}(S_{\alpha'\alpha'}x_{\alpha'})=S_{\alpha''\alpha''}(T_{\alpha''\alpha'}x_{\alpha'}).\\
 35^\alpha. & T_{\alpha''\alpha'}0_{\alpha'}=0_{\alpha''}.
\end{array}
\]
\end{hbtheorem}

\subsection{Definition by primitive recursion: theorems 36--43}

The source formalizes definition by primitive recursion as follows
\cite[\S8, p.~68]{Church1940}: given \(\mathbf A_{\alpha'}\) and
\(\mathbf B_{\alpha'\alpha'\alpha'}\), find \(\mathbf F_{\alpha'\alpha'}\) such
that, with \(\mathbf x_{\alpha'}\) not free in \(\mathbf A_{\alpha'}\),
\(\mathbf B_{\alpha'\alpha'\alpha'}\) or \(\mathbf F_{\alpha'\alpha'}\),
\[
 \derives\mathbf F_{\alpha'\alpha'}0_{\alpha'}=\mathbf A_{\alpha'},\qquad
 \derives N_{o\alpha'}\mathbf x_{\alpha'}\supset
 \mathbf F_{\alpha'\alpha'}(S_{\alpha'\alpha'}\mathbf x_{\alpha'})=
 \mathbf B_{\alpha'\alpha'\alpha'}\mathbf x_{\alpha'}(\mathbf F_{\alpha'\alpha'}\mathbf x_{\alpha'}).
\]

\begin{hbdefinition}{def:p02-w02-p02-040-church-recursor}{Church's primitive-recursion formula}
For \(\mathbf A_{\alpha'}\) and \(\mathbf B_{\alpha'\alpha'\alpha'}\) in which
\(x_{\alpha'}\) and \(y_{\alpha''}\) are not free, the source takes
\(\mathbf F_{\alpha'\alpha'}\) to be \cite[\S8, p.~68]{Church1940}
\[
\begin{aligned}
 \lambda x_{\alpha'}\,.\,T_{\alpha'''\alpha''}(T_{\alpha''\alpha'}x_{\alpha'})
 (\lambda y_{\alpha''}\langle S_{\alpha'\alpha'}(y_{\alpha''}(K_{\alpha'\alpha'\alpha'}I_{\alpha'})0_{\alpha'}),\,
 &\mathbf B_{\alpha'\alpha'\alpha'}(y_{\alpha''}(K_{\alpha'\alpha'\alpha'}I_{\alpha'})0_{\alpha'})
 (y_{\alpha''}(K_{\alpha'\alpha'\alpha'}0_{\alpha'})I_{\alpha'})\rangle)\\
 &\langle0_{\alpha'},\mathbf A_{\alpha'}\rangle(K_{\alpha'\alpha'\alpha'}0_{\alpha'})I_{\alpha'}.
\end{aligned}
\]
Here \(I_{\alpha'}\) is the identity abbreviation \(I_{\gamma\gamma}\) with
\(\gamma=\alpha\alpha\), a wff of type \(\alpha'\).  A pair
\(\langle\mathbf M_{\alpha'},\mathbf N_{\alpha'}\rangle\) has type \(\alpha''\);
applied to \(K_{\alpha'\alpha'\alpha'}I_{\alpha'}\) and \(0_{\alpha'}\) it yields
its first component, and applied to \(K_{\alpha'\alpha'\alpha'}0_{\alpha'}\) and
\(I_{\alpha'}\) its second (theorems 37 and 38).  The formula iterates, a number
of times read off from \(x_{\alpha'}\) at type \(\alpha'''\), the step that
sends a pair (counter, value) to (successor of counter, \(\mathbf B\) of
both), starting from \(\langle0_{\alpha'},\mathbf A_{\alpha'}\rangle\), and
returns the second component.  The predecessor \(P_{\alpha'\alpha'}\) of
\cref{def:p02-w02-p02-040-church-abbrev} is the case
\(\mathbf A_{\alpha'}=0_{\alpha'}\), \(\mathbf B_{\alpha'\alpha'\alpha'}=K_{\alpha'\alpha'\alpha'}\),
whose step returns the old counter \cite[\S8, p.~68]{Church1940}.
\index{primitive recursion!in Church's type theory}
\end{hbdefinition}

\begin{hbtheorem}{thm:p02-w02-p02-040-church-36-43}{proved=annex:p02-w02-p02-040-church-36-43}{Theorems 36--43}
For every type \(\alpha\), using Axioms 1--6\(^\alpha\) only, the following are
theorems \cite[\S8, p.~68]{Church1940}:
\[
\begin{array}{rl}
 36^\alpha. & N_{o\alpha'}n_{\alpha'}\supset\lambda f_{\alpha\alpha}\lambda x_\alpha(n_{\alpha'}f_{\alpha\alpha}x_\alpha)=n_{\alpha'}.\\
 37^\alpha. & N_{o\alpha'}m_{\alpha'}\supset.\,N_{o\alpha'}n_{\alpha'}\supset\langle m_{\alpha'},n_{\alpha'}\rangle(K_{\alpha'\alpha'\alpha'}I_{\alpha'})0_{\alpha'}=m_{\alpha'}.\\
 38^\alpha. & N_{o\alpha'}m_{\alpha'}\supset.\,N_{o\alpha'}n_{\alpha'}\supset\langle m_{\alpha'},n_{\alpha'}\rangle(K_{\alpha'\alpha'\alpha'}0_{\alpha'})I_{\alpha'}=n_{\alpha'}.\\
 39^\alpha. & N_{o\alpha'''}n_{\alpha'''}\supset n_{\alpha'''}(\lambda p_{\alpha''}\langle S_{\alpha'\alpha'}(p_{\alpha''}(K_{\alpha'\alpha'\alpha'}I_{\alpha'})0_{\alpha'}),p_{\alpha''}(K_{\alpha'\alpha'\alpha'}I_{\alpha'})0_{\alpha'}\rangle)\\
   & \qquad\langle0_{\alpha'},0_{\alpha'}\rangle(K_{\alpha'\alpha'\alpha'}I_{\alpha'})0_{\alpha'}
   =n_{\alpha'''}S_{\alpha''\alpha''}0_{\alpha''}S_{\alpha'\alpha'}0_{\alpha'}.\\
 40^\alpha. & N_{o\alpha'''}n_{\alpha'''}\supset P_{\alpha'\alpha'''}(S_{\alpha'''\alpha'''}n_{\alpha'''})=n_{\alpha'''}S_{\alpha''\alpha''}0_{\alpha''}S_{\alpha'\alpha'}0_{\alpha'}.\\
 41^\alpha. & P_{\alpha'\alpha'''}0_{\alpha'''}=0_{\alpha'}.
\end{array}
\]
If \(\derives30^\alpha\) and \(\derives30^{\alpha'}\), then, using in addition
Axioms 9\(^{\alpha''}\) and 9\(^{\alpha'''}\), also
\[
\begin{array}{rl}
 42^\alpha. & P_{\alpha'\alpha'}0_{\alpha'}=0_{\alpha'}.\\
 43^\alpha. & N_{o\alpha'}n_{\alpha'}\supset P_{\alpha'\alpha'}(S_{\alpha'\alpha'}n_{\alpha'})=n_{\alpha'}.
\end{array}
\]
\end{hbtheorem}

The source states 36--43 with a parenthetical list of what each uses and no
further proof; the annex writes each derivation out.  Theorem~36 is the
\(\eta\)-equation restricted to natural numbers; it needs no extensionality
axiom because every number is, up to conversion, an abstract
\(\lambda f_{\alpha\alpha}\lambda x_\alpha(\cdots)\).

\begin{hbtheorem}{thm:p02-w02-p02-040-church-recursion}{proved=annex:p02-w02-p02-040-church-recursion}{The recursion equations for Church's formula}
Let \(\alpha\) be a type with \(\derives30^\alpha\) and \(\derives30^{\alpha'}\),
let \(\mathbf F_{\alpha'\alpha'}\) be the formula of
\cref{def:p02-w02-p02-040-church-recursor}, and let \(\mathbf x_{\alpha'}\) be
a variable not free in \(\mathbf A_{\alpha'}\) or \(\mathbf B_{\alpha'\alpha'\alpha'}\).
Suppose
\[
 \text{(H)}\qquad\derives N_{o\alpha'}\mathbf A_{\alpha'}\quad\text{and}\quad
 \derives N_{o\alpha'}u_{\alpha'}\supset.\,N_{o\alpha'}v_{\alpha'}\supset
 N_{o\alpha'}(\mathbf B_{\alpha'\alpha'\alpha'}u_{\alpha'}v_{\alpha'}),
\]
with \(u_{\alpha'},v_{\alpha'}\) not free in \(\mathbf B_{\alpha'\alpha'\alpha'}\).
Then, using Axioms 1--6\(^\alpha\), 9\(^{\alpha''}\) and 9\(^{\alpha'''}\),
\[
 \derives\mathbf F_{\alpha'\alpha'}0_{\alpha'}=\mathbf A_{\alpha'},\qquad
 \derives N_{o\alpha'}\mathbf x_{\alpha'}\supset
 \mathbf F_{\alpha'\alpha'}(S_{\alpha'\alpha'}\mathbf x_{\alpha'})=
 \mathbf B_{\alpha'\alpha'\alpha'}\mathbf x_{\alpha'}(\mathbf F_{\alpha'\alpha'}\mathbf x_{\alpha'}).
\]
\end{hbtheorem}

\begin{hbwarning}{def:p02-w02-p02-040-church-recursion-hypothesis}{The recursion schema needs a hypothesis}
The source asserts the two recursion equations for arbitrary
\(\mathbf A_{\alpha'}\) and \(\mathbf B_{\alpha'\alpha'\alpha'}\)
\cite[\S8, p.~68]{Church1940}.  For its formula they can fail without a
hypothesis such as (H).  The first projection
\(\langle\mathbf M,\mathbf N\rangle(K_{\alpha'\alpha'\alpha'}I_{\alpha'})0_{\alpha'}\)
converts to \(\lambda h_{\alpha\alpha}\lambda x_\alpha(\mathbf M h_{\alpha\alpha}(\mathbf N I_{\alpha\alpha}x_\alpha))\),
which returns \(\mathbf M\) only when \(\mathbf N I_{\alpha\alpha}x_\alpha\) behaves
like \(x_\alpha\); a value of \(\mathbf B\) that is not a number can corrupt the
counter.  The annex gives a standard model of
\cref{def:p02-w02-p02-040-church-standard} in which, for
\(\alpha=\iota\), \(\mathbf A_{\iota'}\) a closed formula denoting a constant
functional and \(\mathbf B_{\iota'\iota'\iota'}=K_{\iota'\iota'\iota'}\), the
second equation is false at \(\mathbf x_{\iota'}=0_{\iota'}\).  Since every
theorem is valid there, that instance is not a theorem.
\end{hbwarning}

\begin{hbtheorem}{thm:p02-w02-p02-040-church-rosser}{proved=annex:p02-w02-p02-040-church-rosser}{Rosser's addition and multiplication}
The source's n.~13 observes that the recursion formula involves
descriptions through \(T_{\alpha'''\alpha''}\) and \(T_{\alpha''\alpha'}\), and
that for addition and multiplication of non-negative integers the following
definitions due to J.~B.~Rosser avoid them \cite[\S8, p.~68, n.~13]{Church1940}:
\[
\begin{array}{r@{\ \longrightarrow\ }l}
 \mathfrak s_{\alpha'\alpha'\alpha'} & \lambda m_{\alpha'}\lambda n_{\alpha'}\lambda f_{\alpha\alpha}\lambda x_\alpha(m_{\alpha'}f_{\alpha\alpha}(n_{\alpha'}f_{\alpha\alpha}x_\alpha)),\\
 \mathfrak B_{\alpha'\alpha'\alpha'} & \lambda m_{\alpha'}\lambda n_{\alpha'}\lambda f_{\alpha\alpha}(m_{\alpha'}(n_{\alpha'}f_{\alpha\alpha})),\\
 {}[\mathbf A_{\alpha'}+\mathbf B_{\alpha'}] & \mathfrak s_{\alpha'\alpha'\alpha'}\mathbf A_{\alpha'}\mathbf B_{\alpha'},\\
 {}[\mathbf A_{\alpha'}\times\mathbf B_{\alpha'}] & \mathfrak B_{\alpha'\alpha'\alpha'}\mathbf A_{\alpha'}\mathbf B_{\alpha'}.
\end{array}
\]
For every type \(\alpha\), using Axioms 1--6\(^\alpha\) only, the following are
theorems:
\[
\begin{array}{ll}
 N_{o\alpha'}n_{\alpha'}\supset0_{\alpha'}+n_{\alpha'}=n_{\alpha'}, &
 S_{\alpha'\alpha'}m_{\alpha'}+n_{\alpha'}=S_{\alpha'\alpha'}(m_{\alpha'}+n_{\alpha'}),\\
 0_{\alpha'}\times n_{\alpha'}=0_{\alpha'}, &
 S_{\alpha'\alpha'}m_{\alpha'}\times n_{\alpha'}=n_{\alpha'}+(m_{\alpha'}\times n_{\alpha'}),\\
 N_{o\alpha'}m_{\alpha'}\supset.\,N_{o\alpha'}n_{\alpha'}\supset N_{o\alpha'}(m_{\alpha'}+n_{\alpha'}), &
 N_{o\alpha'}m_{\alpha'}\supset.\,N_{o\alpha'}n_{\alpha'}\supset N_{o\alpha'}(m_{\alpha'}\times n_{\alpha'}).
\end{array}
\]
\end{hbtheorem}

These are recursion equations in the first argument.  The source states the
definitions only; the equations displayed are the handbook's justification
that the abbreviations define addition and multiplication of natural numbers,
and they hold at every type \(\alpha\), with no use of Axioms 7--11.

\subsection{Interpretation, consistency and independence}

\begin{hbdefinition}{def:p02-w02-p02-040-church-standard}{Standard interpretation of Church's system}
Let \(D\) be a nonempty set.  Define sets \(\mathcal D_\gamma\) by recursion
on type symbols: \(\mathcal D_\iota=D\), \(\mathcal D_o=\{\mathrm T,\mathrm F\}\)
with \(\mathrm T\ne\mathrm F\), and \(\mathcal D_{\alpha\beta}\) the set of all
functions from \(\mathcal D_\beta\) to \(\mathcal D_\alpha\).  Each
\(\mathcal D_\gamma\) is nonempty.  Choose, for every type \(\alpha\), a
function \(c_\alpha:\mathcal D_{o\alpha}\to\mathcal D_\alpha\) with
\(g(c_\alpha(g))=\mathrm T\) whenever \(g(a)=\mathrm T\) for some
\(a\in\mathcal D_\alpha\).  The \emph{standard model} \(\mathcal S(D,c)\)
interprets \(N_{oo}\) as Boolean negation, \(A_{ooo}\) as the curried Boolean
disjunction (\(A_{ooo}uv=\mathrm T\) iff \(u=\mathrm T\) or \(v=\mathrm T\)),
\(\Pi_{o(o\alpha)}\) as the function sending \(g\in\mathcal D_{o\alpha}\) to
\(\mathrm T\) iff \(g(a)=\mathrm T\) for all \(a\in\mathcal D_\alpha\), and
\(\iota_{\alpha(o\alpha)}\) as \(c_\alpha\).  An assignment \(s\) maps each
variable \(\mathbf x_\gamma\) into \(\mathcal D_\gamma\).  Denotation
\(\sem{\mathbf M}_s\) is defined by
\[
 \sem{\mathbf x}_s=s(\mathbf x),\qquad
 \sem{(\mathbf F\mathbf A)}_s=\sem{\mathbf F}_s(\sem{\mathbf A}_s),\qquad
 \sem{(\lambda\mathbf x_\beta\mathbf M)}_s=
 \bigl(a\mapsto\sem{\mathbf M}_{s[\mathbf x_\beta\mapsto a]}\bigr),
\]
with constants as above.  A wff \(\mathbf M_o\) is \emph{valid} in
\(\mathcal S(D,c)\) iff \(\sem{\mathbf M_o}_s=\mathrm T\) for every \(s\).
\index{simple theory of types!standard model}
\end{hbdefinition}

The choice of the family \((c_\alpha)\) uses the Axiom of Choice in the
metatheory: one choice function on the nonempty subsets of each
\(\mathcal D_\alpha\).  This formalizes the source's intended
interpretation; the source itself proves no soundness or consistency
theorem and states its non-theoremhood claims with ``consistency assumed''
\cite[\S6, p.~64]{Church1940}.

\begin{hbtheorem}{thm:p02-w02-p02-040-church-soundness}{proved=annex:p02-w02-p02-040-church-soundness}{Soundness, consistency and independence}
Work in ZFC.
\begin{enumerate}
 \item In every standard model \(\mathcal S(D,c)\),
   \([\mathbf A_\alpha=\mathbf B_\alpha]\) denotes \(\mathrm T\) under \(s\) iff
   \(\sem{\mathbf A_\alpha}_s=\sem{\mathbf B_\alpha}_s\); Axioms 1--6\(^\alpha\),
   9\(^\alpha\), 10\(^{\alpha\beta}\) and 11\(^\alpha\) are valid; Axiom~7 is
   valid iff \(D\) has at least two members; Axiom~8 is valid if \(D\) is
   infinite or has exactly one member, and is not valid if \(D\) is finite
   with at least two members; and Rules I--VI carry formulas valid in
   \(\mathcal S(D,c)\) to formulas valid in it.
 \item Consequently, every theorem is valid in \(\mathcal S(\N,c)\), and
   the system is consistent: the variable \(p\) is not a theorem.
 \item Axiom~7 is not a theorem of the system whose axioms are all the
   others, and Axiom~8 is not a theorem of the system whose axioms are all
   the others \cite[\S4, p.~61]{Church1940}.
 \item If \(\alpha\) is a type symbol in which \(\iota\) does not occur, then
   \[
    26^\alpha.\quad N_{o\alpha'}x_{\alpha'}\supset.\,N_{o\alpha'}y_{\alpha'}\supset.\,
    S_{\alpha'\alpha'}x_{\alpha'}=S_{\alpha'\alpha'}y_{\alpha'}\supset x_{\alpha'}=y_{\alpha'}
   \]
   is not a theorem \cite[\S6, p.~64]{Church1940}.
\end{enumerate}
\end{hbtheorem}

The source states clauses 3 and 4: the independence of Axiom~7 by an
interpretation with exactly one individual, that of Axiom~8 by one with
finitely many individuals, more than one, and the non-theoremhood of
26\(^\alpha\) for types consisting entirely of \(o\)'s
\cite[\S4, p.~61; \S6, p.~64]{Church1940}.  The annex carries out these
interpretations in standard models.

\subsection{Descriptions and the selection constants}

\begin{hbtheorem}{thm:p02-w02-p02-040-church-9-from-11}{proved=annex:p02-w02-p02-040-church-9-from-11}{Descriptions from choice}
For every type \(\alpha\), 9\(^\alpha\) is derivable from Axioms 1--4 and
11\(^\alpha\) by Rules I, IV\('\) and V.
\end{hbtheorem}

\begin{hbdefinition}{def:p02-w02-p02-040-church-reduced}{The system with two primitive selection constants}
The \emph{reduced language} has the primitive symbols of
\cref{def:p02-w02-p02-040-church-wff} except that the only selection
constants are \(\iota_{o(oo)}\) and \(\iota_{\iota(o\iota)}\).  For every other
type \(\gamma\), which has the form \(\alpha\beta\), the expression
\(\iota_{\alpha\beta(o(\alpha\beta))}\) abbreviates the formula of the schema
of p.~62 displayed above, in which \(\iota_{\alpha(o\alpha)}\) (inside the
description) is primitive or is itself an abbreviation of lower rank; by
recursion on type symbols every \(\iota_{\gamma(o\gamma)}\) thus abbreviates a
closed wff of the reduced language of type \(\gamma(o\gamma)\).  The
\emph{reduced system} has as axioms 1--8, 9\(^o\), 9\(^\iota\) and every
10\(^{\alpha\beta}\), read in the reduced language, and Rules I--VI.
\end{hbdefinition}

\begin{hbtheorem}{thm:p02-w02-p02-040-church-selection}{proved=annex:p02-w02-p02-040-church-selection}{The defined selection constants satisfy 9; independence of 9}
\begin{enumerate}
 \item In the reduced system every formula 9\(^\gamma\), read through the
   abbreviations of \cref{def:p02-w02-p02-040-church-reduced}, is a theorem.
   Its derivation from 9\(^\alpha\) to 9\(^{\alpha\beta}\) uses Axioms
   1--6\(^\gamma\) and 10\(^{\alpha\beta}\) only.
 \item In the full system, for every type \(\alpha\), Axiom 9\(^\alpha\) is not
   derivable from the remaining axioms among 1--8, 9\(^\beta\) (\(\beta\neq\alpha\))
   and 10\(^{\beta\gamma}\); in particular each 9\(^\alpha\) is independent when
   1--9 are the only axioms \cite[\S4, p.~62]{Church1940}.
\end{enumerate}
\end{hbtheorem}

Clause~1 is the claim of the source on p.~62 with a proof; the proof uses
Axiom~10 at every compound type, which is why the source makes the
reduction conditional on the extensionality axioms.  Clause~2 also shows the
sense of the source's remark that the 9\(^\alpha\) for compound \(\alpha\),
``although independent, are superfluous'' once 10 is present: they cannot be
dropped as axioms about primitive constants, but they become theorems when
the constants are defined.

\subsection{Bridge to the host simple-type syntax}

\begin{hbdefinition}{def:p02-w02-p02-040-church-bridge}{Translation of Church's notation into the host syntax}
Define a map \(\gamma\mapsto\gamma^\circ\) from Church's type symbols to the
host simple types of \cref{sec:simple-type-theory-01} by
\[
 \iota^\circ=\iota,\qquad o^\circ=o,\qquad
 (\alpha\beta)^\circ=\beta^\circ\to\alpha^\circ.
\]
The order is reversed: Church writes the value type first, the host writes
the argument type first.  Thus \(o\iota\iota\) becomes
\(\iota\to\iota\to o\) (with \(\to\) associating to the right), and
\(\alpha'\) becomes \((\alpha^\circ\to\alpha^\circ)\to(\alpha^\circ\to\alpha^\circ)\).
Extend the host signature by one nonlogical constant
\(\varepsilon_\alpha:(\alpha^\circ\to o)\to\alpha^\circ\) for each Church type
\(\alpha\).  Translate wffs by
\[
\begin{aligned}
 (\mathbf x_\alpha)^\circ&=\mathbf x{:}\alpha^\circ,
 &(\lambda\mathbf x_\beta\mathbf M)^\circ&=\lambda\mathbf x{:}\beta^\circ.\,\mathbf M^\circ,
 &(\mathbf F\mathbf A)^\circ&=\mathbf F^\circ\mathbf A^\circ,\\
 N_{oo}^\circ&=\neg,
 &A_{ooo}^\circ&=\lor,
 &\Pi_{o(o\alpha)}^\circ&=\forall_{\alpha^\circ},
 \qquad\iota_{\alpha(o\alpha)}^\circ=\varepsilon_\alpha,
\end{aligned}
\]
with \(\neg,\lor\) the constants of \cref{def:p02-stt-boolean-constants} and
\(\forall_{\alpha^\circ}\) the host quantifier constant.
\index{simple theory of types!Church-to-host bridge}
\end{hbdefinition}

\begin{hbtheorem}{thm:p02-w02-p02-040-church-bridge}{proved=annex:p02-w02-p02-040-church-bridge}{The bridge preserves types and denotations}
The type map is a bijection between Church's type symbols and the host
simple types.  For every Church wff \(\mathbf M\) of type \(\gamma\),
\(\mathbf M^\circ\) is a well-typed host term of type \(\gamma^\circ\), and
free and bound variables correspond.  Let \(\mathcal M\) be a host standard
type structure with \(D_\iota=D\) in which each \(\varepsilon_\alpha\)
denotes a function \(c_\alpha\) as in
\cref{def:p02-w02-p02-040-church-standard}, the truth values \(1,0\) of the
host playing the roles of \(\mathrm T,\mathrm F\).  Then for every
assignment \(s\),
\(\sem{\mathbf M^\circ}^{\mathcal M}_s=\sem{\mathbf M}_s\) computed in
\(\mathcal S(D,c)\).  In particular, the translation of Church's defined
equality \([\mathbf A_\alpha=\mathbf B_\alpha]\) has the same denotation as
the host term \(\mathbf A^\circ=_{\alpha^\circ}\mathbf B^\circ\), and the
translations of Church's defined \(\supset\), conjunction and \(\exists\)
have the same denotations as the host \(\to\), \(\land\) and
\(\exists_{\alpha^\circ}\).
\end{hbtheorem}

\begin{hbwarning}{def:p02-w02-p02-040-church-bridge-limits}{What the bridge does not identify}
Three differences survive the translation.  First, the host lists
\(=_\alpha\), \(\bottom\), \(\land\) and \(\to\) as primitive constants,
whereas Church defines equality by the Leibniz formula \(Q_{o\alpha\alpha}\)
and defines \(\supset\) and conjunction from \(N_{oo}\) and \(A_{ooo}\); the
denotations agree in standard structures by
\cref{thm:p02-w02-p02-040-church-soundness}.  They also agree in every host
Henkin type structure, but for a host-specific reason: there the closure
clause puts the denotation of \(\lambda y{:}\alpha.\,x=_\alpha y\) into
\(D_{\alpha\to o}\), and instantiating the Leibniz formula at that predicate
yields identity.  Church's language has no primitive identity, so the
same argument is not available for a selected-domain reading of his
system; the general models of \cref{def:p02-w02-p02-040-church-general}
impose the separation condition (G5) instead, which makes Leibniz equality
identity (\cref{thm:p02-w02-p02-040-church-general-soundness}), and
\cref{thm:p02-w02-p02-040-church-general-completeness} shows that every
consistent set has such a model.  Second, the host treats \(\eta\)-conversion as a
typed operation, whereas Church's conversion rules are \(\alpha\) and
\(\beta\) only and the \(\eta\)-equation is the theorem
21\(^{\alpha\beta}\), which depends on Axiom 10\(^{\alpha\beta}\).  Third,
Church's selection constant is logical and governed by Axioms 9 and 11,
while its image \(\varepsilon_\alpha\) is a nonlogical host constant: a host
standard structure may interpret it by a function that is not a choice
function, and then the translation of 11\(^\alpha\) is false there.
\end{hbwarning}

\begin{hbexample}{def:p02-w02-p02-040-church-bridge-example}{Translating the choice axiom and the numerals}
Axiom 11\(^\alpha\), \(f_{o\alpha}x_\alpha\supset f_{o\alpha}(\iota_{\alpha(o\alpha)}f_{o\alpha})\),
abbreviates \(A_{ooo}(N_{oo}(f_{o\alpha}x_\alpha))(f_{o\alpha}(\iota_{\alpha(o\alpha)}f_{o\alpha}))\).
With \(\alpha=\iota\) its translation is
\[
 \lor\bigl(\neg(f\,x)\bigr)\bigl(f(\varepsilon_\iota f)\bigr),
 \qquad f:\iota\to o,\ x:\iota,
\]
which is true under every assignment exactly when \(\varepsilon_\iota\)
denotes a choice function on the nonempty sets of individuals.  The numeral
\(2_{\iota'}\) is \(\lambda f_{\iota\iota}\lambda x_\iota(f_{\iota\iota}(f_{\iota\iota}x_\iota))\);
its translation
\(\lambda f{:}\iota\to\iota.\,\lambda x{:}\iota.\,f(f\,x)\) has type
\((\iota\to\iota)\to(\iota\to\iota)\), the image of \(\iota'\).  The
successor \(S_{\iota'\iota'}\) translates to
\(\lambda n.\,\lambda f.\,\lambda x.\,f(n\,f\,x)\) of type
\(\iota'^\circ\to\iota'^\circ\).
\end{hbexample}

\begin{hbexample}{def:p02-w02-p02-040-church-bridge-arith}{The arithmetic theorems in host notation}
Write \(\nu=\alpha^\circ\), so that \(\alpha'^\circ=(\nu\to\nu)\to(\nu\to\nu)\).
The translation of \(N_{o\alpha'}\) is the host term
\(\lambda n.\,\forall_{\alpha'^\circ\to o}(\lambda f.\,\lor(\neg(f\,\bar0))(\lor(\neg\forall(\lambda x.\,\lor(\neg(f\,x))(f(\mathsf s\,x))))(f\,n)))\),
with \(\bar0\) and \(\mathsf s\) the translations of \(0_{\alpha'}\) and
\(S_{\alpha'\alpha'}\); in a host standard structure it denotes the set of
Church numerals \(\varphi\mapsto\varphi^n\) over \(D_\nu\), by
\cref{thm:p02-w02-p02-040-church-bridge} and the numeral clause in the proof
of \cref{thm:p02-w02-p02-040-church-soundness}.  Hence the translation of
26\(^\alpha\) says that \(\mathsf s\) is injective on these numerals, that of
30\(^\alpha\) that \(m\mapsto m\,\mathsf s\,\bar0\) is injective on the numerals
of type \(\alpha''^\circ\), that of 32\(^\alpha\) that the translated
\(T_{\alpha''\alpha'}\) inverts this map on numerals, and those of 43\(^\alpha\)
and of Rosser's equations that the translated \(P_{\alpha'\alpha'}\),
\(\mathfrak s\) and \(\mathfrak B\) act on numerals as predecessor, addition
and multiplication.  Because every theorem is valid in \(\mathcal S(\N,c)\),
each statement of \cref{thm:p02-w02-p02-040-church-26-30,thm:p02-w02-p02-040-church-31-35,thm:p02-w02-p02-040-church-36-43,thm:p02-w02-p02-040-church-rosser}
holds, under this reading, in every host standard structure with
\(D_\iota=\N\) whose constants \(\varepsilon_\gamma\) denote choice functions,
for the types \(\alpha\) at which it was derived.  The bridge transfers truth
in standard structures, not derivability: the host fixes no calculus for
which these translations are claimed to be theorems.
\end{hbexample}

\begin{hbwarning}{def:p02-w02-p02-040-church-boundaries}{Language boundaries of the source}
The object language has no type variables: a superscripted schema such as
5\(^\alpha\) stands for infinitely many axioms, and a proof of a schema is
a uniform recipe for proofs of its instances, the device the source calls
typical ambiguity \cite[\S5, p.~63, n.~10]{Church1940}.  Types are not
cumulative: the natural numbers of type \(\alpha'\) and of type
\(\alpha''\) are different objects, connected by \(T_{\alpha''\alpha'}\).
A complete incorporation of the calculus of \(\lambda\)-conversion is
impossible while \(\lambda\) and juxtaposition keep their meanings, so the
incorporation is partial \cite[p.~56]{Church1940}.  Which types satisfy
26\(^\alpha\) is delicate: it is a theorem for \(\alpha=\iota,\iota',\iota'',\ldots\)
(\cref{thm:p02-w02-p02-040-church-26-30}) and not for types built from \(o\)
alone (\cref{thm:p02-w02-p02-040-church-soundness}), and theorems 31--35 and
42--43 are therefore available only at types for which 30 has been proved.
\end{hbwarning}

\subsection{The system read as a typed higher-order calculus}

Read as a higher-order calculus, the system has the following components,
each displayed above.  The language is that of
\cref{def:p02-w02-p02-040-church-wff}, with variables of every type and so
quantification over functions and propositional functions of every finite
type.  Derivability is the theorem relation of
\cref{def:p02-w02-p02-040-church-axioms}, and derivability from assumptions
is the restricted relation defined there, for which the deduction theorem
holds.  Conversion is \(\alpha\)- and \(\beta\)-conversion, Rules I--III.
Comprehension is not an axiom: it is the theorem
\cref{thm:p02-w02-p02-040-church-comprehension}, obtained from abstraction,
conversion and 13, and it holds for every wff, quantifiers of every type
included.  Choice is Axiom 11\(^\alpha\), a selection operator at every type;
descriptions, Axiom 9\(^\alpha\), are its special case for propositional
functions true of exactly one object.  Extensionality is Axiom
10\(^{\alpha\beta}\) for functions; the \(\eta\)-equation is its consequence
21\(^{\alpha\beta}\); and extensionality for propositions is not assumed.
Equality is the defined Leibniz relation, whose reflexivity, substitutivity,
congruence, symmetry and transitivity are 16--20.  Axioms 7 and 8 assert
infinity.

The relational calculus of \cref{def:p02-sol-derived-calculus}, with
the schema families of \cref{def:p02-sol-matched-henkin}, differs on
each of these points: it has relation variables only, no abstraction, a
comprehension schema postulated outright, graph choice in place of a
selection operator, and primitive object equality.  The standard model of
\cref{def:p02-w02-p02-040-church-standard} is the full semantics of the
typed language: every function type ranges over all set-theoretic
functions.  The general models of the next subsection are the matching
selected-domain semantics, and the calculus is sound and complete for them
(\cref{thm:p02-w02-p02-040-church-general-soundness,thm:p02-w02-p02-040-church-general-completeness}).

\subsection{General models, the first-order encoding and completeness}
\label{sec:p02-w02-p02-040-church-general}

\emph{Purpose.}  The standard models of
\cref{def:p02-w02-p02-040-church-standard} let every function type range
over all functions.  This subsection widens the class of models so that
function types range over selected sets of functions, closed enough that
every wff denotes, and proves that Church's calculus, exactly as displayed
in \cref{def:p02-w02-p02-040-church-rules,def:p02-w02-p02-040-church-axioms},
is sound and complete for the wider class.  The completeness proof passes
through an exact many-sorted first-order encoding: general models and
models of one first-order theory are carried into each other by an encoding
and a decoding that preserve truth.
\emph{Inputs.}  The rules and axioms of Church's system; the deduction
theorem (\cref{thm:p02-w02-p02-040-church-deduction}); strict many-sorted
first-order logic (\cref{def:p02-many-sorted-signature}); ZF set theory as
metatheory.
\emph{Conclusions.}  Soundness
(\cref{thm:p02-w02-p02-040-church-general-soundness}); the encoding and
decoding theorem (\cref{thm:p02-w02-p02-040-church-encoding}); and strong
completeness with compactness
(\cref{thm:p02-w02-p02-040-church-general-completeness}).
\emph{Standing.}  The source contains no model theory of its system; the
definitions, theorems and proofs of this subsection are the handbook's.

\begin{hbdefinition}{def:p02-w02-p02-040-church-general}{General models of Church's system}
A \emph{general model} \(\mathcal M\) consists of the following data.
\begin{enumerate}
 \item[(G1)] For every type symbol \(\gamma\) a nonempty set
   \(\mathcal D_\gamma\), such that every member of \(\mathcal D_{\alpha\beta}\)
   is a function from \(\mathcal D_\beta\) to \(\mathcal D_\alpha\).
 \item[(G2)] A set \(\mathcal T\subseteq\mathcal D_o\), the \emph{true}
   propositions.
 \item[(G3)] Members \(n\in\mathcal D_{oo}\), \(a\in\mathcal D_{ooo}\),
   \(\pi_\alpha\in\mathcal D_{o(o\alpha)}\) and
   \(c_\alpha\in\mathcal D_{\alpha(o\alpha)}\) for every \(\alpha\), such that
   for all \(u,v\in\mathcal D_o\) and \(g\in\mathcal D_{o\alpha}\):
   \(n(u)\in\mathcal T\) iff \(u\notin\mathcal T\);
   \(a(u)(v)\in\mathcal T\) iff \(u\in\mathcal T\) or \(v\in\mathcal T\);
   \(\pi_\alpha(g)\in\mathcal T\) iff \(g(e)\in\mathcal T\) for every
   \(e\in\mathcal D_\alpha\); and \(g(c_\alpha(g))\in\mathcal T\) whenever
   \(g(e)\in\mathcal T\) for some \(e\in\mathcal D_\alpha\).
 \item[(G4)] (\emph{Closure.})  For every wff \(\mathbf M_\gamma\) and every
   assignment \(s\), that is, every map sending each variable
   \(\mathbf x_\delta\) into \(\mathcal D_\delta\), the three denotation
   clauses of \cref{def:p02-w02-p02-040-church-standard}, with
   \(N_{oo}\), \(A_{ooo}\), \(\Pi_{o(o\alpha)}\) and \(\iota_{\alpha(o\alpha)}\)
   denoting \(n\), \(a\), \(\pi_\alpha\) and \(c_\alpha\), define a member
   \(\sem{\mathbf M_\gamma}_s\) of \(\mathcal D_\gamma\): at every abstraction
   \((\lambda\mathbf x_\beta\mathbf M_\alpha)\) the function
   \(e\mapsto\sem{\mathbf M_\alpha}_{s[\mathbf x_\beta\mapsto e]}\) belongs to
   \(\mathcal D_{\alpha\beta}\).
 \item[(G5)] (\emph{Separation.})  For every \(\alpha\) and all distinct
   \(e,e'\in\mathcal D_\alpha\) there is \(h\in\mathcal D_{o\alpha}\) with
   \(h(e)\in\mathcal T\) and \(h(e')\notin\mathcal T\).
\end{enumerate}
Write \(\mathcal M,s\entails\mathbf A_o\) when
\(\sem{\mathbf A_o}_s\in\mathcal T\); \(\mathbf A_o\) is \emph{valid} in
\(\mathcal M\) when this holds for every assignment \(s\).  For a set
\(\Gamma\) of wffs of type \(o\), write \(\Gamma\entails_{\mathrm G}\mathbf B_o\)
when for every general model \(\mathcal M\) in which Axioms 7 and 8 are valid
and every assignment \(s\) with \(\mathcal M,s\entails\mathbf A_o\) for all
\(\mathbf A_o\in\Gamma\), also \(\mathcal M,s\entails\mathbf B_o\).
\index{simple theory of types!general model}
\index{Henkin semantics!for Church's type theory}
\end{hbdefinition}

Every standard model \(\mathcal S(D,c)\) is a general model with
\(\mathcal T=\{\mathrm T\}\): (G4) holds because each
\(\mathcal D_{\alpha\beta}\) contains every function, and (G5) because the
characteristic function of \(\{e\}\) lies in \(\mathcal D_{o\alpha}\).  Three
features of the definition carry the requirements of the calculus.
\emph{Nonempty domains} (G1) are needed for Axiom~5\(^\alpha\) and Rule~IV,
which instantiate a variable; in the encoding below they become the nonempty
sorts of many-sorted first-order logic.  \emph{Extensionality} for
functions, Axiom~10\(^{\alpha\beta}\), is matched by two things together:
members of \(\mathcal D_{\alpha\beta}\) are functions, so two of them that
agree everywhere are identical, and (G5) makes Church's defined Leibniz
equality denote identity.  \emph{Extensionality for propositions} is not
imposed: \(\mathcal D_o\) may have more than two members, several of them
true.  If \(\mathcal D_o\) has exactly two members, \(p\equiv q\supset p=q\)
is valid, and the source treats that formula as a possible further axiom,
not as one of its axioms \cite[\S4, p.~61]{Church1940}.

\begin{hbdefinition}{def:p02-w02-p02-040-church-encoding}{The first-order encoding of Church's system}
Fix an enumeration of the variables of Church's language.  The strict
many-sorted signature \(\Sigma_{\mathrm C}\) (in the sense of
\cref{def:p02-many-sorted-signature}) has one sort \(\mathsf s_\gamma\) for each type
symbol \(\gamma\) and the following symbols:
\begin{itemize}
 \item for all \(\alpha,\beta\) a function symbol
   \(\mathsf{ap}_{\alpha\beta}:\mathsf s_{\alpha\beta}\times\mathsf s_\beta\to\mathsf s_\alpha\),
   written \(\mathsf{ap}\) when the sorts are evident;
 \item constants \(\widehat N:\mathsf s_{oo}\), \(\widehat A:\mathsf s_{ooo}\),
   \(\widehat\Pi_\alpha:\mathsf s_{o(o\alpha)}\) and
   \(\widehat\iota_\alpha:\mathsf s_{\alpha(o\alpha)}\) for every \(\alpha\);
 \item for every wff \(\mathbf L=(\lambda\mathbf x_\beta\mathbf M_\alpha)\) of the
   primitive notation, with free variables
   \(\mathbf y^1_{\delta_1},\ldots,\mathbf y^k_{\delta_k}\) listed in the order of
   the enumeration, a function symbol
   \(\ell_{\mathbf L}:\mathsf s_{\delta_1}\times\cdots\times\mathsf s_{\delta_k}\to\mathsf s_{\alpha\beta}\),
   a constant when \(k=0\);
 \item a unary predicate symbol \(\mathsf{True}\) on \(\mathsf s_o\).
\end{itemize}
Each Church variable \(\mathbf x_\gamma\) also serves as a variable of sort
\(\mathsf s_\gamma\).  The \emph{encoding} \(\mathbf M\mapsto\mathbf M^\dagger\),
applied after all abbreviations are eliminated, sends wffs to
\(\Sigma_{\mathrm C}\)-terms:
\[
 \mathbf x^\dagger=\mathbf x,\quad
 N_{oo}^\dagger=\widehat N,\quad A_{ooo}^\dagger=\widehat A,\quad
 \Pi_{o(o\alpha)}^\dagger=\widehat\Pi_\alpha,\quad
 \iota_{\alpha(o\alpha)}^\dagger=\widehat\iota_\alpha,\quad
 (\mathbf F\mathbf A)^\dagger=\mathsf{ap}(\mathbf F^\dagger,\mathbf A^\dagger),\quad
 \mathbf L^\dagger=\ell_{\mathbf L}(\mathbf y^1,\ldots,\mathbf y^k).
\]
Thus \(\mathbf M_\gamma^\dagger\) has sort \(\mathsf s_\gamma\) and the same free
variables as \(\mathbf M_\gamma\).  For each \(\gamma\) let
\[
 \mathsf{Lb}_\gamma(u,v)\defeq\forall h{:}\mathsf s_{o\gamma}\,
 \bigl(\mathsf{True}(\mathsf{ap}(h,u))\to\mathsf{True}(\mathsf{ap}(h,v))\bigr).
\]
The theory \(\mathsf{Enc}\) consists of the universal closures of the
following formulas, for all type symbols \(\alpha,\beta\) and all abstracts
\(\mathbf L=(\lambda\mathbf x_\beta\mathbf M_\alpha)\):
\[
\begin{array}{rl}
 (\mathrm E_N) & \mathsf{True}(\mathsf{ap}(\widehat N,p))\leftrightarrow\neg\mathsf{True}(p),\\
 (\mathrm E_A) & \mathsf{True}(\mathsf{ap}(\mathsf{ap}(\widehat A,p),q))\leftrightarrow(\mathsf{True}(p)\lor\mathsf{True}(q)),\\
 (\mathrm E_\Pi^\alpha) & \mathsf{True}(\mathsf{ap}(\widehat\Pi_\alpha,f))\leftrightarrow\forall x{:}\mathsf s_\alpha\,\mathsf{True}(\mathsf{ap}(f,x)),\\
 (\mathrm E_\iota^\alpha) & \mathsf{True}(\mathsf{ap}(f,x))\to\mathsf{True}(\mathsf{ap}(f,\mathsf{ap}(\widehat\iota_\alpha,f))),\\
 (\mathrm E_\beta^{\mathbf L}) & \mathsf{ap}(\ell_{\mathbf L}(\mathbf y^1,\ldots,\mathbf y^k),\mathbf x_\beta)=\mathbf M_\alpha^\dagger,\\
 (\mathrm E_{\mathrm{ext}}^{\alpha\beta}) & \forall x{:}\mathsf s_\beta\,\mathsf{Lb}_\alpha(\mathsf{ap}(f,x),\mathsf{ap}(g,x))\to\mathsf{Lb}_{\alpha\beta}(f,g),\\
 (\mathrm E_7),\ (\mathrm E_8) & \mathsf{True}(7^\dagger),\quad\mathsf{True}(8^\dagger),
\end{array}
\]
where \(p,q\) have sort \(\mathsf s_o\), \(x\) sort \(\mathsf s_\alpha\) in
\((\mathrm E_\Pi^\alpha)\) and \((\mathrm E_\iota^\alpha)\), \(f,g\) the evident
function sorts, and \(7^\dagger,8^\dagger\) are the encodings of Axioms 7 and
8.
\index{simple theory of types!first-order encoding}
\end{hbdefinition}

Each component of the encoding has one job.  The symbols \(\ell_{\mathbf L}\)
with \((\mathrm E_\beta^{\mathbf L})\) are the first-order form of closure
(G4): every abstract, with any values of its parameters, names an element,
and applying that element computes the body.  Comprehension
(\cref{thm:p02-w02-p02-040-church-comprehension}) is therefore built into
the signature rather than postulated as a schema.  The axioms
\((\mathrm E_N)\), \((\mathrm E_A)\), \((\mathrm E_\Pi^\alpha)\) are (G3) for the
logical constants, \((\mathrm E_\iota^\alpha)\) is the choice condition, and
\((\mathrm E_{\mathrm{ext}}^{\alpha\beta})\) is extensionality stated with the
first-order Leibniz formula \(\mathsf{Lb}\).  Church's defined equality is
encoded faithfully: in every model of \(\mathsf{Enc}\),
\(\mathsf{True}([\mathbf A_\alpha=\mathbf B_\alpha]^\dagger)\) holds exactly when
\(\mathsf{Lb}_\alpha(\mathbf A^\dagger,\mathbf B^\dagger)\) does (shown in the
proof of \cref{thm:p02-w02-p02-040-church-encoding}).

\begin{hbtheorem}{thm:p02-w02-p02-040-church-general-soundness}{proved=annex:p02-w02-p02-040-church-general-soundness}{Soundness for general models}
In every general model, \([\mathbf A_\alpha=\mathbf B_\alpha]\) is true under
\(s\) iff \(\sem{\mathbf A_\alpha}_s=\sem{\mathbf B_\alpha}_s\), and Axioms
1--6\(^\alpha\), 9\(^\alpha\), 10\(^{\alpha\beta}\) and 11\(^\alpha\) are valid.  If
\(\mathbf A^1_o,\ldots,\mathbf A^n_o\derives\mathbf B_o\), then
\(\{\mathbf A^1_o,\ldots,\mathbf A^n_o\}\entails_{\mathrm G}\mathbf B_o\).  In
particular every theorem is valid in every general model in which Axioms 7
and 8 are valid.
\end{hbtheorem}

\begin{hbtheorem}{thm:p02-w02-p02-040-church-encoding}{proved=annex:p02-w02-p02-040-church-encoding}{Encoding and decoding}
\begin{enumerate}
 \item (Encoding.)  Let \(\mathcal M\) be a general model in which Axioms 7
   and 8 are valid, and fix an assignment \(s_0\).  Let \(\mathcal M^\dagger\)
   interpret \(\mathsf s_\gamma\) as \(\mathcal D_\gamma\), \(\mathsf{ap}\) as
   application, \(\widehat N,\widehat A,\widehat\Pi_\alpha,\widehat\iota_\alpha\) as
   \(n,a,\pi_\alpha,c_\alpha\), \(\ell_{\mathbf L}(e_1,\ldots,e_k)\) as
   \(\sem{\mathbf L}_{s_0[\mathbf y^1\mapsto e_1,\ldots,\mathbf y^k\mapsto e_k]}\),
   \(\mathsf{True}\) as \(\mathcal T\), and \(=\) as identity.  Then
   \(\mathcal M^\dagger\) is a model of \(\mathsf{Enc}\), and for every wff
   \(\mathbf M\) and assignment \(s\) the value of \(\mathbf M^\dagger\) under
   \(s\) is \(\sem{\mathbf M}_s\).  Hence
   \(\mathcal M,s\entails\mathbf A_o\) iff
   \(\mathcal M^\dagger,s\entails\mathsf{True}(\mathbf A_o^\dagger)\).
 \item (Decoding.)  Let \(\mathfrak A\) be a model of \(\mathsf{Enc}\) and
   \(t_0\) an \(\mathfrak A\)-assignment.  There are a general model
   \(\mathfrak A^\natural\), in which Axioms 7 and 8 are valid, and
   surjections \(d_\gamma\) from the domain of \(\mathsf s_\gamma\) in
   \(\mathfrak A\) onto the domain \(\mathcal D_\gamma\) of
   \(\mathfrak A^\natural\), such that \(d_\gamma(u)=d_\gamma(v)\) iff
   \(\mathfrak A\entails\mathsf{Lb}_\gamma(u,v)\), and such that for every wff
   \(\mathbf M\) and every \(\mathfrak A\)-assignment \(t\),
   \(\sem{\mathbf M}_{d\circ t}\) is \(d\) applied to the value of
   \(\mathbf M^\dagger\) under \(t\).  Hence
   \(\mathfrak A,t\entails\mathsf{True}(\mathbf A_o^\dagger)\) iff
   \(\mathfrak A^\natural,d\circ t\entails\mathbf A_o\).
 \item (Exactness.)  For \(\mathcal M\) as in clause~1, the maps
   \(d_\gamma\) of clause~2 for \(\mathfrak A=\mathcal M^\dagger\) are
   bijections that commute with application and carry \(n,a,\pi_\alpha,c_\alpha\)
   and \(\mathcal T\) to their counterparts, so the decoding of
   \(\mathcal M^\dagger\) is isomorphic to \(\mathcal M\).
\end{enumerate}
\end{hbtheorem}

\begin{hbtheorem}{thm:p02-w02-p02-040-church-general-completeness}{proved=annex:p02-w02-p02-040-church-general-completeness}{Strong completeness for general models}
For every set \(\Gamma\) of wffs of type \(o\) and every wff \(\mathbf B_o\),
the following are equivalent.
\begin{enumerate}
 \item \(\mathbf A^1_o,\ldots,\mathbf A^n_o\derives\mathbf B_o\) for some
   \(n\geq0\) and some \(\mathbf A^1_o,\ldots,\mathbf A^n_o\in\Gamma\).
 \item \(\Gamma\entails_{\mathrm G}\mathbf B_o\).
 \item For every model \(\mathfrak A\) of \(\mathsf{Enc}\) and every
   \(\mathfrak A\)-assignment \(t\), if
   \(\mathfrak A,t\entails\mathsf{True}(\mathbf A_o^\dagger)\) for every
   \(\mathbf A_o\in\Gamma\), then
   \(\mathfrak A,t\entails\mathsf{True}(\mathbf B_o^\dagger)\).
\end{enumerate}
The proof that 3 implies 1 uses no form of the Axiom of Choice: the
countermodel is built from the wffs of the language.  Consequently the
theorems of Church's system are exactly the wffs of type \(o\) valid in every
general model in which Axioms 7 and 8 are valid, and \(\entails_{\mathrm G}\)
is compact.
\end{hbtheorem}

The relation in clause~1 is Church's derivability on assumptions, with its
restriction that no variable free in an assumption is substituted for or
generalized upon; the semantic relation \(\entails_{\mathrm G}\) matches it by
holding the values of those variables fixed in one assignment \(s\).

\begin{hbexample}{def:p02-w02-p02-040-church-general-example}{Encoding choice and one comprehension instance}
Axiom 11\(^\iota\) abbreviates
\(A_{ooo}(N_{oo}(f_{o\iota}x_\iota))(f_{o\iota}(\iota_{\iota(o\iota)}f_{o\iota}))\).
Its encoding is
\(\mathsf{True}(\mathsf{ap}(\mathsf{ap}(\widehat A,\mathsf{ap}(\widehat N,\mathsf{ap}(f,x))),\mathsf{ap}(f,\mathsf{ap}(\widehat\iota_\iota,f))))\),
which by \((\mathrm E_A)\) and \((\mathrm E_N)\) is equivalent in every model of
\(\mathsf{Enc}\) to
\(\neg\mathsf{True}(\mathsf{ap}(f,x))\lor\mathsf{True}(\mathsf{ap}(f,\mathsf{ap}(\widehat\iota_\iota,f)))\),
the axiom \((\mathrm E_\iota^\iota)\).  For comprehension, take
\(\mathbf L=\lambda x_\iota(N_{oo}(g_{o\iota}x_\iota))\), whose only free
variable is \(g_{o\iota}\).  Then \(\mathbf L^\dagger=\ell_{\mathbf L}(g)\), and
\((\mathrm E_\beta^{\mathbf L})\) is
\(\mathsf{ap}(\ell_{\mathbf L}(g),x)=\mathsf{ap}(\widehat N,\mathsf{ap}(g,x))\).
With \((\mathrm E_N)\) it gives
\(\forall g\,\forall x\,(\mathsf{True}(\mathsf{ap}(\ell_{\mathbf L}(g),x))\leftrightarrow\neg\mathsf{True}(\mathsf{ap}(g,x)))\):
every propositional function of individuals has a complement among the
elements of sort \(\mathsf s_{o\iota}\).  In the decoding, \(\ell_{\mathbf L}(g)\)
becomes the member of \(\mathcal D_{o\iota}\) true exactly where the decoded
\(g\) is not, the witness of the comprehension instance
\((\exists f_{o\iota})(x_\iota)[f_{o\iota}x_\iota\equiv\,\sim g_{o\iota}x_\iota]\).
\end{hbexample}

\begin{hbwarning}{def:p02-w02-p02-040-church-general-limits}{What general-model completeness does not give}
First, the restriction on proofs on assumptions cannot be dropped: from the
assumption \(p\) the formula \((p)p\) does not follow.  In the standard model
\(\mathcal S(\N,c)\), a general model in which Axioms 7 and 8 are valid
(\cref{thm:p02-w02-p02-040-church-soundness}), an assignment with
\(s(p)=\mathrm T\) makes \(p\) true and \((p)p\) false, so
\(\{p\}\not\entails_{\mathrm G}(p)p\).  Second, the theorem characterizes the
theorems by truth in all general models; it says nothing about which wffs
are true in all standard models, where the domains are full function spaces.
Third, \(\mathsf{Enc}\) has one function symbol and one axiom
\((\mathrm E_\beta^{\mathbf L})\) for each of infinitely many abstracts; no
finite part of it encodes the whole system.
\end{hbwarning}

\section{The relational second-order sheet as a native calculus}
\label{sec:p02-w02-p02-040-sol-native}

\emph{Purpose.}  \Cref{def:p02-sol-derived-calculus} defines derivability for
the relational second-order language by translation into many-sorted
first-order logic, and \cref{thm:p02-sol-derived-adequacy} proves that
relation sound and strongly complete for matched Henkin models; its
boundary statement \cref{def:p02-sol-native-boundary} records that the
rules displayed on \texttt{sheet:sol}, used as an
independent generator list, need a converse deductive correspondence.
This section proves that correspondence by proving the displayed rules
complete on their own.
\emph{Inputs.}  The relational language of
\cref{sec:second-order-logic-01} over an at most countable nonlogical
signature; the initial sequents and rules of
\texttt{sheet:lk-eq}; relation substitution
(\cref{def:p02-sol-relation-substitution}); matched Henkin models
(\cref{def:p02-sol-matched-henkin}).
\emph{Conclusions.}  Native soundness and completeness for finite sequents,
without the Axiom of Choice
(\cref{thm:p02-w02-p02-040-sol-native}); and, with the Axiom of Choice as in
\cref{thm:p02-sol-derived-adequacy}, coincidence of native and
translation-defined derivability
(\cref{thm:p02-w02-p02-040-sol-correspondence}).

\begin{hbdefinition}{def:p02-w02-p02-040-sol-native}{Native derivability on the relational sheet}
A sequent \(\Sigma\Rightarrow\Pi\), with \(\Sigma\) and \(\Pi\) finite multisets
of relational second-order formulas, is \emph{natively derivable} if it is
the root of a finite tree whose leaves are initial sequents of
\texttt{sheet:lk-eq} or sequents
\(\Rightarrow\sigma\) with \(\sigma\) an eligible instance of Comp or Choice
from \texttt{sheet:sol}, and each of whose other nodes
follows from its children by a rule of \texttt{sheet:lk-eq} or by one of
\(\forall^2L\), \(\forall^2R\), \(\exists^2L\), \(\exists^2R\).  The rules of
\texttt{sheet:lk-eq} are read over relational second-order formulas: their
formula letters range over all such formulas, their term letters over
object terms, their eigenvariables are object variables, and \([t/x]\) is
capture-avoiding substitution of an object term for an object variable,
which renames object binders when needed and passes through relation
binders unchanged.  For a set \(\Gamma\) of formulas,
\(\Gamma\derives_{\mathrm{nat}}\varphi\) means that \(\Sigma\Rightarrow\varphi\) is
natively derivable for some finite multiset \(\Sigma\) of members of
\(\Gamma\).  A sequent is \emph{\(\mathsf H_{\mathsf{CC}}\)-valid} if every
\(\mathsf H_{\mathsf{CC}}\)-model and matching assignment that satisfy every
member of \(\Sigma\) satisfy some member of \(\Pi\).
\end{hbdefinition}

\begin{hbtheorem}{thm:p02-w02-p02-040-sol-native}{proved=annex:p02-w02-p02-040-sol-native}{Soundness and completeness of the native relational calculus}
Let the nonlogical signature be at most countable.  A sequent
\(\Sigma\Rightarrow\Pi\) is natively derivable iff it is
\(\mathsf H_{\mathsf{CC}}\)-valid.  The proof uses no form of the Axiom of
Choice.
\end{hbtheorem}

\begin{hbcorollary}{thm:p02-w02-p02-040-sol-correspondence}{proved=annex:p02-w02-p02-040-sol-correspondence}{Native and translation-defined derivability coincide}
Assume the Axiom of Choice in the ambient set theory, and let the
nonlogical signature be at most countable.  For every set \(\Gamma\) of
formulas and every formula \(\varphi\),
\[
 \Gamma\derives_{\mathrm{nat}}\varphi
 \quad\Longleftrightarrow\quad
 \Gamma\derives_{\mathsf{SOL}_{\mathrm H}^{\mathsf{CC}}}\varphi
 \quad\Longleftrightarrow\quad
 \Gamma\entails_{\mathrm H}^{\mathsf{CC}}\varphi ,
\]
and a finite sequent \(\Sigma\Rightarrow\Pi\) is natively derivable iff it is
derivable on the linked sheet in the translation-defined sense of
\cref{def:p02-sol-derived-calculus}.
\end{hbcorollary}

So the four relation-quantifier rules, Comp and Choice, added to
\texttt{sheet:lk-eq}, generate exactly the translation-defined relation; the
converse deductive correspondence that
\cref{def:p02-sol-native-boundary} names as missing is the first equivalence
of the corollary.

\begin{hbexample}{def:p02-w02-p02-040-sol-native-example}{A native derivation that needs Comp}
Read \(\varphi\leftrightarrow\psi\) as \((\varphi\to\psi)\land(\psi\to\varphi)\).
The sequent \(\Rightarrow\exists X^1\forall x\,X^1(x)\) is natively derivable:
\[
\begin{array}{rll}
 1. & \Rightarrow y=y,\ Z(y) & {=}I,\ W_R\\
 2. & Z(y)\Rightarrow Z(y) & \mathrm{id}\\
 3. & y=y\to Z(y)\Rightarrow Z(y) & {\to}L\ \text{from 1, 2}\\
 4. & (Z(y)\to y=y)\land(y=y\to Z(y))\Rightarrow Z(y) & W_L,\ {\land}L\\
 5. & \forall x\,(Z(x)\leftrightarrow x=x)\Rightarrow Z(y) & \forall L,\ t:=y\\
 6. & \forall x\,(Z(x)\leftrightarrow x=x)\Rightarrow\forall x\,Z(x) & \forall R,\ \text{eigenvariable }y\\
 7. & \forall x\,(Z(x)\leftrightarrow x=x)\Rightarrow\exists X^1\forall x\,X^1(x) & \exists^2R,\ U:=Z\\
 8. & \exists X^1\forall x\,(X^1(x)\leftrightarrow x=x)\Rightarrow\exists X^1\forall x\,X^1(x) & \exists^2L,\ \text{eigenvariable }Z\\
 9. & \Rightarrow\exists X^1\forall x\,(X^1(x)\leftrightarrow x=x) & \text{Comp},\ \chi(x):=x=x\\
 10. & \Rightarrow\exists X^1\forall x\,X^1(x) & \text{Cut, 9 and 8}
\end{array}
\]
Here \(Z\) is a unary relation variable and \(y\) an object variable; the
eigenvariable conditions hold because \(y\) is not free in line~6's
other formulas and \(Z\) is not free in line~8.  The Comp leaf is
essential: a derivation whose leaves are all initial sequents of
\texttt{sheet:lk-eq} preserves truth in every Henkin structure of
\cref{sec:second-order-logic-05}, since the soundness argument in the proof
of \cref{thm:p02-w02-p02-040-sol-native} uses the matched closure conditions
only at Comp and Choice leaves, and the Henkin structure with \(M=\{0\}\)
and \(\mathcal R_n=\{\varnothing\}\) for every \(n\) makes
\(\exists X^1\forall x\,X^1(x)\) false.  So every derivation of this sequent
has a Comp or Choice leaf.
\end{hbexample}

\begin{hbwarning}{def:p02-w02-p02-040-sol-native-limits}{Limits of the native completeness theorem}
The theorem is completeness for matched Henkin models, not for full
semantics: no sound effective calculus is complete for full second-order
validity (\cref{thm:p02-w02-p02-040-no-effective-calculus}).  The completeness
proof uses Cut, and Cut cannot be dropped from the equality calculus even at
first order: \cref{thm:p07-eqlk-cut-failure} gives a valid first-order
sequent \(a=b\Rightarrow b=a\) with no Cut/Mix-free derivation on
\texttt{sheet:lk-eq}.  Its backward invariant (only the atoms \(a=b\),
\(a=a\) and \(b=a\) occur) survives the added rules, because the four
relation-quantifier rules have a quantified principal formula in their
conclusion and every Comp or Choice leaf contains a quantified formula; so
that sequent has no Cut/Mix-free native derivation either.
The countability hypothesis is used to enumerate the
formulas without choice; the corollary inherits the Axiom of Choice only
from \cref{thm:p02-sol-derived-adequacy}.
\end{hbwarning}

\section{Full second-order categoricity with addition, multiplication and completeness}
\label{sec:p02-w02-p02-040-categoricity}

This section proves the two categoricity theorems used in
\cref{sec:second-order-logic-04}: the arithmetic one for the full signature
\(\{0,S,+,\times\}\) of \cref{def:p02-pa2}, and the one for the second-order
ordered field.  \emph{Inputs.}  Full semantics
(\cref{sec:second-order-logic-02}); the successor categoricity proof of
\cref{sec:second-order-logic-03}; ZF set theory with the natural numbers,
integers and rationals as metatheory.  \emph{Conclusions.}
\Cref{thm:p02-w02-p02-040-pa2-categorical,thm:p02-w02-p02-040-reals-categorical}.
Both conclusions depend on full semantics: the proofs apply a second-order
axiom to subsets of the domain that a Henkin structure need not contain.

\begin{hbtheorem}{thm:p02-w02-p02-040-pa2-categorical}{proved=annex:p02-w02-p02-040-pa2-categorical}{Full categoricity of \(\mathsf{PA}_2\) with addition and multiplication}
The structure \(\mathcal N=(\N,0,n\mapsto n+1,+,\times)\) satisfies
\(\mathsf{PA}_2\) of \cref{def:p02-pa2} under full semantics.  If
\(\mathcal M=(M,0^{\mathcal M},S^{\mathcal M},+^{\mathcal M},\times^{\mathcal M})\)
satisfies \(\mathsf{PA}_2\) under full semantics, the map
\(f:\N\to M\) with \(f(0)=0^{\mathcal M}\) and
\(f(n+1)=S^{\mathcal M}(f(n))\) is an isomorphism from \(\mathcal N\) onto
\(\mathcal M\): it is a bijection preserving \(0\) and \(S\), and for all
\(m,n\in\N\),
\[
 f(m+n)=f(m)+^{\mathcal M}f(n),\qquad
 f(m\times n)=f(m)\times^{\mathcal M}f(n).
\]
Hence any two full models of \(\mathsf{PA}_2\) are isomorphic.
\end{hbtheorem}

\begin{hbdefinition}{def:p02-w02-p02-040-ordered-field}{The ordered-field axioms}
In the signature \(\Lang_{\mathsf R}=\{0,1,+,\times,<\}\) write \(xy\) for
\(x\times y\).  The \emph{ordered-field axioms} are the universal closures
of the following fifteen first-order formulas:
\[
\begin{array}{ll}
 \text{(F1)}\ x+(y+z)=(x+y)+z, & \text{(F2)}\ x+y=y+x,\\
 \text{(F3)}\ x+0=x, & \text{(F4)}\ \exists y\,(x+y=0),\\
 \text{(F5)}\ x(yz)=(xy)z, & \text{(F6)}\ xy=yx,\\
 \text{(F7)}\ x1=x, & \text{(F8)}\ x\neq0\to\exists y\,(xy=1),\\
 \text{(F9)}\ x(y+z)=xy+xz, & \text{(F10)}\ 0\neq1,\\
 \text{(O1)}\ \neg(x<x), & \text{(O2)}\ (x<y\land y<z)\to x<z,\\
 \text{(O3)}\ x<y\lor x=y\lor y<x, & \text{(O4)}\ x<y\to x+z<y+z,\\
 \text{(O5)}\ (0<x\land0<y)\to0<xy. &
\end{array}
\]
The \emph{full second-order ordered-field theory} adds the completeness
sentence displayed in \cref{thm:second-order-reals-categorical}, with
\(u\le v\) abbreviating \(u<v\lor u=v\).  A structure satisfying these
sixteen sentences under full semantics is a \emph{complete ordered field}.
\index{ordered field!axioms}
\index{complete ordered field!second-order}
\end{hbdefinition}

\begin{hbtheorem}{thm:p02-w02-p02-040-reals-categorical}{proved=annex:p02-w02-p02-040-reals-categorical}{Full categoricity and uncountability of complete ordered fields}
Let \(\mathcal F\) and \(\mathcal G\) be complete ordered fields in the
sense of \cref{def:p02-w02-p02-040-ordered-field}.  For \(x\in F\) put
\(L_{\mathcal F}(x)=\{q\in\mathbb Q:q_{\mathcal F}<x\}\), where
\(q\mapsto q_{\mathcal F}\) is the embedding of the rationals into
\(\mathcal F\) constructed in the annex.  Then
\[
 h(x)=\sup\nolimits_{\mathcal G}\{q_{\mathcal G}:q\in L_{\mathcal F}(x)\}
\]
defines an isomorphism of \(\Lang_{\mathsf R}\)-structures from
\(\mathcal F\) onto \(\mathcal G\), the unique map with
\(L_{\mathcal G}(h(x))=L_{\mathcal F}(x)\) for every \(x\).  Moreover the
domain of every complete ordered field is uncountable.
\end{hbtheorem}

This theorem supplies the proof of
\cref{thm:second-order-reals-categorical}: the map is the cut map described
there, the axioms abbreviated there are those of
\cref{def:p02-w02-p02-040-ordered-field}, and the uncountability is proved
directly from the completeness sentence.  The clause there identifying
every model with \(\mathbb R\) presupposes that the ambient real field is a
complete ordered field; given any one complete ordered field, the theorem
makes every other isomorphic to it.

\section{The nonaxiomatizability of full second-order validity}
\label{sec:p02-w02-p02-040-nonaxiomatizable}

\emph{Inputs.}  \Cref{thm:p02-w02-p02-040-pa2-categorical}; the G\"odel
numbering of \cref{sec:arithmetic-coding-04}; computably enumerable sets
as domains of partial computable functions (\cref{sec:turing-machines-03});
\cref{thm:p04-ce-sigma-one}; and \cref{thm:tarski-undefinability}.
\emph{Conclusion.}  \Cref{thm:p02-w02-p02-040-no-effective-calculus}, which
is the theorem stated in \cref{sec:second-order-logic-04} with its
reduction and effectiveness proved.

\begin{hbdefinition}{def:p02-w02-p02-040-so-coding}{Codes of second-order arithmetic formulas}
Extend the tree coding of \cref{sec:arithmetic-coding-04}, with its
constructor \(N_k(a_0,\ldots,a_{r-1})=C(k,[a_0,\ldots,a_{r-1}])\), to the
relational second-order language over \(\{0,S,+,\times\}\) by four new
tags: the relation variable \(X^n_i\) has code \(N_{15}(n,i)\); the atom
\(X^n_i(t_1,\ldots,t_n)\) has code
\(N_{16}(N_{15}(n,i),[\ulcorner t_1\urcorner,\ldots,\ulcorner t_n\urcorner])\);
and \(\forall X^n_i\varphi\), \(\exists X^n_i\varphi\) have codes
\(N_{17}(n,i,\ulcorner\varphi\urcorner)\), \(N_{18}(n,i,\ulcorner\varphi\urcorner)\).
All first-order terms and formulas keep the codes given in
\cref{sec:arithmetic-coding-04}.  Write
\(\mathrm{Val}_2=\{\ulcorner\varphi\urcorner:\varphi\text{ a second-order sentence and }\entails_{\mathrm{full}}\varphi\}\)
and \(\mathrm{TA}=\{\ulcorner\sigma\urcorner:\sigma\text{ a first-order }\Lang_A\text{-sentence and }\mathcal N\entails\sigma\}\).
\end{hbdefinition}

\begin{hbtheorem}{thm:p02-w02-p02-040-no-effective-calculus}{proved=annex:p02-w02-p02-040-no-effective-calculus}{No effective complete calculus for full second-order validity}
\begin{enumerate}
 \item The function \(r(c)=N_{10}(\ulcorner\mathsf{PA}_2\urcorner,c)\), which
   sends the code of \(\sigma\) to the code of \(\mathsf{PA}_2\to\sigma\), is
   primitive recursive, and the set of codes of first-order
   \(\Lang_A\)-sentences is decidable.
 \item For every first-order \(\Lang_A\)-sentence \(\sigma\),
   \(\entails_{\mathrm{full}}(\mathsf{PA}_2\to\sigma)\) iff
   \(\mathcal N\entails\sigma\).
 \item \(\mathrm{TA}\) is not computably enumerable.
 \item \(\mathrm{Val}_2\) is not computably enumerable.
 \item There is no calculus \(\mathsf D\) whose set of theorems (as codes) is
   computably enumerable and which satisfies
   \(\mathsf D\derives\sigma\Leftrightarrow\entails_{\mathrm{full}}\sigma\) for
   every second-order sentence \(\sigma\) of this signature.
\end{enumerate}
\end{hbtheorem}

The proof of clause 3 runs through \cref{thm:p04-ce-sigma-one}: a computably
enumerable set is \(\Sigma^0_1\)-definable in \(\mathcal N\), so a
computably enumerable \(\mathrm{TA}\) would have an arithmetical truth
definition, which \cref{thm:tarski-undefinability} excludes; this is the
corollary that follows that theorem.  Clause 5 concerns effectiveness only:
the set of all full validities, taken as axioms with no rules, is a
complete calculus whose theorem set is not computably enumerable.
Henkin-type completeness, as in \cref{thm:p02-sol-derived-adequacy}, is
completeness for a different consequence relation and is not affected.

\begin{hbexample}{def:p02-w02-p02-040-reduction-example}{One instance of the reduction}
For \(\sigma\defeq\forall x\,\forall y\,(x+y=y+x)\), clause 2 of
\cref{thm:p02-w02-p02-040-no-effective-calculus} gives
\(\entails_{\mathrm{full}}(\mathsf{PA}_2\to\sigma)\), because addition on
\(\N\) is commutative.  A single instance shows only how the reduction
acts; the separation comes from clause 3: by \cref{thm:tarski-undefinability} no
computably enumerable set of first-order sentences contains exactly the
true ones, so for every sound effective calculus some true \(\sigma\) gives
a full validity \(\mathsf{PA}_2\to\sigma\) that the calculus does not derive.
\end{hbexample}
